% concatenated from main.tex, experiments_summary_story.tex, appendix_convergence_proofs.tex, appendix_dca.tex, appendix_prox_linear_stationarity.tex, experiments_appendix_v2.tex, experiments/figures/tab_aggregate.tex, experiments/figures/tab_lts_aggregate.tex, experiments/figures/tab_relu_aggregate.tex
\documentclass{article}
% \newif\ifneuripsfinal
% \neuripsfinalfalse
% \newif\ifneuripspreprint
% \neuripspreprintfalse
% \ifneuripsfinal
%   \usepackage[final]{Formatting_Instructions_For_NeurIPS_2026/neurips_2026}
% \else
%   \ifneuripspreprint
\usepackage[preprint]{Formatting_Instructions_For_NeurIPS_2026/neurips_2026}
%   \else
%     \usepackage{Formatting_Instructions_For_NeurIPS_2026/neurips_2026}
%   \fi
% \fi
\usepackage{amsmath,amsthm,amssymb}
\usepackage{graphicx}
\usepackage{booktabs}
\usepackage{tabularx}
\usepackage{algorithm}
\usepackage{xcolor}
\usepackage{algpseudocode}
\usepackage{hyperref}
\usepackage{cleveref}
\usepackage{import}
\newtheorem{theorem}{Theorem}[section]
\newtheorem{lemma}[theorem]{Lemma}
\newtheorem{definition}{Definition}[section]
\newtheorem{assumption}{Assumption}
\newtheorem{remark}{Remark}[section]
\newcommand{\R}{\mathbb{R}}
\newcommand{\N}{\mathbb{N}}
\DeclareMathOperator{\ri}{ri}
\DeclareMathOperator{\cl}{cl}
\DeclareMathOperator{\aff}{aff}
\DeclareMathOperator{\conv}{conv}
\newcommand{\norm}[1]{\left\|#1\right\|}
\newcommand{\inner}[2]{\langle #1,\, #2 \rangle}
\newcommand{\partialC}{\partial^C}
\newcommand{\partialL}{\partial^L}
\algnewcommand\algorithmicComment[1]{\textit{// #1}}
\algnewcommand{\COMMENT}[1]{\State \algorithmicComment{#1}}
\title{Achieving Directional-Stationarity from a Single Random Direction Step}
\author{Dan Greenstein \qquad Nadav Hallak\\[4pt]
Faculty of Data and Decision Sciences, Technion\\
\texttt{\{sdngreen, ndvhllk\}@technion.ac.il}}
\date{}
\begin{document}
\maketitle
\begin{abstract}
This paper addresses the challenge of obtaining strong optimality guarantees in constrained nonsmooth nonconvex optimization under mild regularity conditions, namely local Lipschitz continuity and existence and continuity of directional derivatives. 
While standard methods typically ensure weak stationarity notions, achieving directional (d-)stationarity remains nontrivial.
We show that a random direction exploration step is sufficient to attain d-stationarity. 
The proposed approach augments any base optimization method with a single exploration step that samples a direction and step size and accepts the candidate based on a function value comparison.
The resulting scheme guarantees that all accumulation points are d-stationary almost surely, independently of the behavior of the underlying method. 
Moreover, it preserves convergence rates of the base method, as established for DCA and prox-linear-type schemes.
The theoretical results are complemented by numerical experiments illustrating the effect and guarantees of the exploration step.
\end{abstract}




\section{Introduction}\label{sec:intro}
\subsection{Problem formulation}
We consider the constrained nonsmooth nonconvex optimization problem
\begin{equation}\label{eq:main}
  \min_{x \in C} h(x),
\end{equation}
where $h \colon \mathbb{R}^n \to \mathbb{R}$ is locally Lipschitz and directionally differentiable, and $C \subseteq \mathbb{R}^n$ is a nonempty closed convex set. 
This setting covers a broad class of problems, including smooth-minus-convex objectives $h = f - g$ with $f \in C^1$ and $g$ convex, as well as the difference-of-convex (DC) class where both $f$ and $g$ are convex. 
Importantly, our approach does not require access to such a decomposition and applies directly to the objective $h$.

\medskip

A central challenge in nonsmooth nonconvex optimization is to obtain strong and meaningful optimality guarantees. In DC programming and related settings, several stationarity notions are available, with varying levels of strength (see, e.g., \citep{Pang2017}). 
Among these, \emph{directional (d-)stationarity} characterizes points with no feasible descent directions, defined through directional derivatives (Definition~\ref{def:d-stationary} and \Cref{ass:A2}).

Directional stationarity is stronger than the criticality notion standard in DC algorithms \citep{deOliveira2020} and than Clarke stationarity \citep{clarke1990optimization}; derivative-free schemes such as MADS~\citep{Audet2006} typically target the latter. See Remark~\ref{rem:stationarity-comparison} for limiting (Mordukhovich) stationarity and for a fuller discussion.
It is rarely guaranteed in general constrained nonsmooth nonconvex settings.
Existing methods achieving d-stationarity either rely on explicit DC decompositions and access to smooth component functions \citep{Pang2017, Liu2019a}, or are restricted to linear feasibility sets \citep{Beck2020, Beck2022}, and their guarantees are tightly coupled to the algorithmic dynamics.

In contrast, we propose a minimal exploration step that is independent of the base algorithm and enforces d-stationarity of accumulation points. 
Hence, any method can be made asymptotically d-stationary without modifying the oracle map itself; the augmentation modifies only the accepted iterate.


\paragraph{Main contributions.}

\begin{itemize}
\item \textbf{Exploration-augmented optimization.}
We introduce a one-step random exploration mechanism that can be integrated into any feasible optimization method.

\item \textbf{D-stationarity via random exploration.}
We prove that, almost surely, all accumulation points are d-stationary, independently of the underlying algorithm. 
The result %relies on a probabilistic argument based on Borel--Cantelli--L\'{e}vy and 
does not use any structural properties of the underlying method.

\item \textbf{Rate preservation and oracle compatibility.}
We establish that the augmentation can preserve the convergence behavior of the base method and apply across different oracle classes such as DC and prox-linear-type schemes.
\end{itemize}


\subsection{Mathematical preliminaries}
We formally define two assumptions on \eqref{eq:main} and one supporting lemma that ensure well-defined directional derivatives and the continuity required for the probabilistic argument.
\begin{assumption}[A1]\label{ass:A1}
$h$ is locally Lipschitz and bounded below on $C$.
\end{assumption}
\begin{assumption}[A2]\label{ass:A2}
The directional derivative $h'(x;\, d) = \lim_{t \downarrow 0}(h(x + td) - h(x))/t$ exists for every $x \in C$ and $d \in \R^n$.
\end{assumption}
\begin{lemma}
    \label{ass:A3}
For every fixed $x \in C$, the map $d \mapsto h'(x;\, d)$ is (locally Lipschitz) continuous.
\end{lemma}
These properties hold for a broad class of objectives such as $h = f - g$ with $f \in C^1$ and $g$ convex and finite-valued on a neighborhood of $C$, as well as in the DC case where both $f$ and $g$ are convex. 
In both settings, the directional derivative admits the representation $h'(x;\, d) = f'(x;\, d) - g'(x;\, d)$, and continuity in $d$ follows from continuity/convexity respectively \citep[Theorems~23.1 and~10.1]{Rockafellar1997}.
% \begin{remark}
% Assumptions~A1--A3 hold for a broad class of objectives such as when $h = f - g$ with $f \in C^1$ and $g$ convex and finite-valued on a neighborhood of $C$, as well as in the DC case where both $f$ and $g$ are convex. In both settings, the directional derivative admits the representation $h'(x;\, d) = \langle \nabla f(x), d \rangle - g'(x;\, d)$ (or $f'(x;\, d) - g'(x;\, d)$ in the DC case), and continuity in $d$ follows from convexity \citep[Theorems~23.1 and~10.1]{Rockafellar1997}.
% \end{remark}
\begin{definition}[Affine hull subspace]
    Given $x^0 \in C$, the linear subspace parallel to $\aff(C)$ is defined by $L := \aff(C - x^0)$. 
\end{definition}
Note that $L$ is independent of the choice of $x^0 \in C$: for any $x^0,x^1 \in C$ the sets $C-x^0$ and $C-x^1$ differ by translation by $x^0-x^1\in \aff(C-x^0)$, hence have the same (linear) affine hull.

 For $x \in C$ and $v \in \R^n$, denote the maximal feasible step length at $x$ in direction $v$ by
\[
  \tau(x;\, v) := \sup\!\bigl\{t \geq 0 : x + tv \in C\bigr\}.
\]

Our analysis uses the cone of feasible directions; a well-known fact is that it is convex.
\begin{definition}[Cone of feasible directions]\label{def:feasible-cone}
Let $C \subseteq \R^n$ be convex and $\bar{x} \in C$. The \emph{cone of feasible directions} is $D_C(\bar{x}) := \{w \mid \exists\, \lambda > 0,\;\bar{x} + \lambda w \in C\}$.
\end{definition}

\begin{definition}[Directional (d-)stationarity]\label{def:d-stationary}
A point $\bar{x} \in C$ is \textit{d-stationary} for \eqref{eq:main} if
\begin{equation*}
  h'(\bar{x};\, d) \geq 0 \qquad \forall\, d \in D_C(\bar{x}).
\end{equation*}
\end{definition}
By Lemma~\ref{ass:A3}, the same condition is equivalent to $h'(\bar{x};\, d) \geq 0$ for every $d \in \cl(D_C(\bar{x}))$; see Lemma~\ref{lem:continuity}.

\begin{remark}[Comparison with other stationarity notions]\label{rem:stationarity-comparison}
For locally Lipschitz $h$ and closed convex $C$, \emph{Clarke stationarity} at $\bar{x} \in C$ is 
$0 \in \partialC h(\bar{x}) + N_C(\bar{x})$ \citep{clarke1990optimization}. 
Writing $h^\circ$ for the Clarke directional derivative, one has $h'(\bar{x}; d) \leq h^\circ(\bar{x}; d)$ for all $d$, with equality if $h$ is Clarke regular, hence Clarke stationarity does not in general imply d-stationarity. 
For convex $C$, d-stationarity is equivalent to $h'(\bar{x}; d) \geq 0$ for all $d \in T_C(\bar{x})$.
\emph{Mordukhovich (limiting) stationarity} ($0 \in \partialL h(\bar{x}) + N_C(\bar{x})$ \citep{RockafellarWets2009}) implies Clarke stationarity but remains weaker than d-stationarity in general.
For DC objectives $h=f-g$ with $f,g$ convex, \emph{criticality} is again weaker than d-stationarity \citep{deOliveira2020}.
\end{remark}

% \begin{remark}[Comparison with other stationarity notions]\label{rem:stationarity-comparison}
% For locally Lipschitz $h$ and closed convex $C$, \textit{Clarke stationarity} at $\bar{x} \in C$ is the condition $0 \in \partialC h(\bar{x}) + N_C(\bar{x})$ in the sense of \citet{clarke1990optimization}. It is a first-order necessary condition for a local constrained minimum \citep[Corollary to Proposition~2.4.3]{clarke1990optimization}.
% Writing $h^\circ(\bar{x};\, \cdot)$ for the Clarke generalized directional derivative, one always has $h'(\bar{x};\, d) \leq h^\circ(\bar{x};\, d)$ for all $d$ \citep[Proposition~2.2.2]{clarke1990optimization}, with equality for every $d$ if and only if $h$ is \textit{regular} at $\bar{x}$ in Clarke's sense \citep[Definition~2.2.4]{clarke1990optimization}. Hence Clarke stationarity does not by itself imply d-stationarity in Definition~\ref{def:d-stationary}. For convex $C$, the tangent cone satisfies $T_C(\bar{x}) = \cl(D_C(\bar{x}))$, so under Assumption~A3 the equivalent formulation on $\cl(D_C(\bar{x}))$ from Lemma~\ref{lem:continuity} is exactly nonnegativity of $h'(\bar{x};\, \cdot)$ on $T_C(\bar{x})$.

% \textit{Mordukhovich} or \textit{limiting} stationarity is the condition $0 \in \partialL h(\bar{x}) + N_C(\bar{x})$ with $\partialL$ the limiting subdifferential \citep{RockafellarWets2009}. For locally Lipschitz $h$, $\partialC h(\bar{x}) = \conv \partialL h(\bar{x})$ \citep[Theorem~8.49]{RockafellarWets2009}, so limiting stationarity implies Clarke stationarity, strictly in general. Theorem~8.15 of \citet{RockafellarWets2009} relates tangent-cone geometry to limiting subgradients; when $h$ is not subdifferentially regular \citep[Definition~7.25]{RockafellarWets2009}, purely subdifferential stationarity can be strictly weaker than requiring $h'(\bar{x};\, d) \geq 0$ for all $d \in T_C(\bar{x})$, as discussed in the commentary opening Chapter~10, Section~A, there.

% For DC objectives $h=f-g$ with $f,g$ convex, the \textit{criticality} condition used in DC algorithms is weaker than d-stationarity; see \citep{deOliveira2020} and Section~\ref{sec:dc} for the subgradient formulation in the constrained DCA setting.
% \end{remark}

Finally, throughout this paper we fix the parameters $\gamma > 0$ and $r \in (0, \infty]$.


\section{Related Work}\label{sec:related}

We summarize the closest strands of the literature in \Cref{tab:comparison}.
To the best of our knowledge, no prior work achieves d-stationarity via random feasible-direction exploration under general convex constraints with oracle-independent guarantees.



\paragraph{DC programming.}
The standard optimality notion in DC algorithms is criticality, which is strictly weaker than d-stationarity \citep{deOliveira2020}. \citet{Pang2017} initiated a line of work establishing d-stationarity for structured DC problems in which the concave part admits a finite-max representation, by solving one convex subproblem per active component. This framework has been extended along several directions by \citet{AhnPangXin2017}, %studied bicriteria sparse learning under the same structure; 
\citet{Liu2019a},% introduced perturbations to reduce the per-iteration complexity to a single subproblem; 
\citet{TaoLi2023}, % obtained linear convergence rates under additional error-bound conditions; 
\citet{SunWu2024}, % proposed hybrid schemes that selectively invoke multi-subproblem steps;  
and \citet{FengYuan2026}, % established almost-sure d-stationarity of cluster points via refined perturbation arguments. 
all of which rely critically on the finite-max structure of the DC decomposition, explicit access to the DC decomposition and its components, and first-order information, and are inherently algorithm-specific rather than oracle-agnostic. 

A complementary line of work \citep{Beck2020, Beck2022} establishes d-stationarity for smooth-minus-convex objectives over polyhedral sets using feasible-direction methods, including randomized variants based on fixed finite spanning sets. In contrast, bundle methods~\citep{Joki2018, deOliveira2021} typically guarantee only criticality or Clarke stationarity.

Finally, while the convergence rates of DCA are well understood in the unconstrained setting \citep{Abbaszadehpeivasti2024}, these results concern criticality measures and do not address d-stationarity. 
In Section~\ref{sec:dc}, we show that the proposed augmentation preserves the $O(1/\sqrt{N})$ rate of DCA in a strongly convex regime, while the d-stationarity upgrade follows from the exploration mechanism and is independent of the DC structure.

% \paragraph{DC programming.}
% The standard optimality notion in DC algorithms is criticality ($\partial f(\bar{x}) \cap \partial g(\bar{x}) \neq \emptyset$), which is strictly weaker than d-stationarity. \citet{Pang2017} initiated a line of work establishing d-stationarity for the structured DC class in which $g$ takes the form $g(x) = \max_{i \in I} \psi_i(x)$ with each $\psi_i \in C^1$, by solving one convex subproblem per active piece. This template has been refined along several complementary axes: \citet{AhnPangXin2017} characterized d-stationarity in bicriteria sparse learning with the same finite-max concave part; \citet{Liu2019a} introduced random perturbation to resolve the active index with a single subproblem per iteration; \citet{TaoLi2023} obtained R-linear rates under piecewise error-bound and isocost conditions; \citet{SunWu2024} proposed hybrid schemes that defer the costly multi-subproblem active-set step until needed; and \citet{FengYuan2026} recently sharpened the perturbation analysis to give d-stationarity of every cluster point almost surely. All of these methods rely critically on the finite-max structure of $g$, on explicit knowledge of the DC decomposition, and on first-order access to the component functions $\psi_i$, and each is a stand-alone algorithm rather than an oracle augmentation. A complementary structural restriction is studied by \citet{Beck2020, Beck2022}: smooth-minus-convex objectives $h = f - g$ with $f \in C^1$ (possibly nonconvex) and $g$ convex (possibly nonsmooth), over polyhedral constraints $\{x : Ax \leq b\}$, where d-stationarity is attained by direct-search feasible-directions methods whose randomized variant polls a fixed finite positive spanning set rather than the full sphere. Bundle methods~\citep{Joki2018, deOliveira2021} target criticality or Clarke stationarity. The unconstrained DCA convergence rate was analyzed in detail by \citet{Abbaszadehpeivasti2024}, who used semidefinite programming performance estimation to establish tight $O(1/\sqrt{N})$ rates on the gradient criticality gap $\min_k \norm{g^k_1 - g^k_2}$ across several combinations of smoothness and strong-convexity assumptions on $f$ and $g$. Section~\ref{sec:dc} treats the constrained case in one such regime, where at least one of the convex parts is strongly convex, and shows that the exploration augmentation preserves a rate of the same $O(1/\sqrt{N})$ order while upgrading cluster points to d-stationarity; we make no claim to matching the constants of \citet{Abbaszadehpeivasti2024} or to covering their full range of regimes.

\paragraph{Direct search: GSS, pattern search, and MADS.}
Generating set search~\citep{Kolda2006} and pattern search~\citep{Lewis2000} poll finitely many directions and yield KKT-type limits in the smooth case and Clarke stationary limits for Lipschitz objectives \citep[Corollary to Proposition~2.4.3]{clarke1990optimization}. MADS~\citep{Audet2006} randomizes to asymptotically dense polling (Clarke stationarity); although the direction coverage is conceptually related to our argument, MADS is a standalone mesh-based method with sufficient decrease, whereas we augment an arbitrary oracle and target d-stationarity under general convex constraints. The feasible-direction methods of \citet{Beck2020, Beck2022} reach d-stationarity only under smooth-minus-convex structure over polyhedral sets.

% \paragraph{Direct search: GSS, pattern search, and MADS.}
% Generating set search~\citep{Kolda2006} and pattern search~\citep{Lewis2000} poll a fixed finite spanning set for descent, establishing subsequence convergence to a KKT point for smooth objectives and to Clarke stationary points for Lipschitz objectives. MADS~\citep{Audet2006} is the randomized extension and the structurally closest prior method to ours: it generates an asymptotically dense set of polling directions in the hypertangent cone, so every interior direction is eventually polled with positive probability, yielding convergence to Clarke stationary points. The mechanism is closely analogous to our Borel--Cantelli argument. However, MADS targets Clarke stationarity (not d-stationarity), is a standalone algorithm (not an oracle augmentation), and uses a shrinking mesh with sufficient decrease rather than a single random step. The methods of \citet{Beck2020, Beck2022}, discussed above, fit methodologically within this family but reach d-stationarity by exploiting their smooth-minus-convex polyhedral structure.

\paragraph{Randomized search outside direct search.}
\citet{Gratton2015} analyze probabilistic-descent direct search, yielding almost-sure convergence and $O(1/\sqrt{k})$ gradient-norm rates for smooth objectives; \citet{Gratton2019} treat linear constraints analogously. These schemes remain direct-search templates targeting gradient- or Clarke-type stationarity. Classical random search \citep{SolisWets1981} and sufficient-descent schemes \citep{Wardi1989} give value- or stationarity-type guarantees under coverage or internal descent, but are standalone algorithms. Our augmentation instead layers one random feasible-direction probe on an arbitrary oracle and targets d-stationarity of accumulation points under convex constraints.

% \paragraph{Probabilistic direct search.}
% \citet{Gratton2015} analyze direct-search methods with probabilistic descent: with positive probability, at least one polling direction yields descent, giving almost-sure convergence and $O(1/\sqrt{k})$ gradient-norm rates for smooth objectives. The constrained extension~\citep{Gratton2019} introduces probabilistic feasible descent for linearly constrained problems.

% \paragraph{Classical random search and sufficient descent.}
% \citet{SolisWets1981} establish asymptotic guarantees in function value and hitting probability of near-optimal regions, under coverage-type assumptions on the sampling process; this is conceptually aligned with our requirement that random exploration not ignore any open subset of feasible directions. In contrast, our target condition is d-stationarity of accumulation points under convex constraints and nonsmooth nonconvex objectives, so the relevant geometric object is the feasible direction set rather than the iterate domain. Sufficient-descent random search methods, such as \citet{Wardi1989}, are closer in spirit to our almost-sure accumulation-point guarantees, but they are stand-alone algorithmic templates whose inner-loop design drives convergence. Our contribution is orthogonal: a wrapper that can be overlaid on a base method without prescribing the base step construction.

\begin{table}[t]
\centering
\caption{Comparison of related fields. %The combination of d-stationarity with a random sphere-direction mechanism for nonsmooth nonconvex objectives under Assumptions~A1--A3 (including DC as a special case) does not appear prior to this work.
}
\label{tab:comparison}
\smallskip
\begin{tabularx}{\linewidth}{@{}>{\raggedright\arraybackslash}X>{\raggedright\arraybackslash}X>{\raggedright\arraybackslash}X>{\raggedright\arraybackslash}X@{}}
\toprule
Field / Work & Target & Mechanism & Nonsmooth DC? \\
\midrule
DC methods \citep{Pang2017} & D-stationarity & Convex subproblem & $g = \max_i \psi_i$, $\psi_i \in C^1$ only \\
GSS / GPS \citep{Kolda2006} & Clarke stat. & Finite spanning set & No \\
MADS \citep{Audet2006} & Clarke stat. & Dense random poll & No \\
Prob.\ direct search \citep{Gratton2015} & Grad.\ stat. & Random sphere dir. & No \\
Random search \citep{SolisWets1981, Wardi1989} & Value / stationarity & Coverage / sufficient descent & No \\
\midrule
This work & D-stationarity & Random sphere dir. & Yes \\
\bottomrule
\end{tabularx}
\end{table}

% \section{Setup and Preliminaries}\label{sec:setup}

% \subsection{Problem and Assumptions}

% As a first step, we formally define the three assumptions  on our model \eqref{eq:main}.

% \begin{assumption}[A1]\label{ass:A1}
% $h$ is locally Lipschitz and bounded below on $C$.
% \end{assumption}

% \begin{assumption}[A2]\label{ass:A2}
% The directional derivative $h'(x;\, d) = \lim_{t \downarrow 0}(h(x + td) - h(x))/t$ exists for every $x \in C$ and $d \in \R^n$.
% \end{assumption}

% \begin{assumption}[A3]\label{ass:A3}
% For every fixed $x \in C$, the map $d \mapsto h'(x;\, d)$ is continuous.
% \end{assumption}

% \begin{remark}
% Assumptions~A1--A3 hold for a broad class of objectives. In particular, they are satisfied when $h = f - g$ with $f \in C^1$ and $g$ convex and finite-valued on a neighborhood of $C$, as well as in the DC case where both $f$ and $g$ are convex. In both settings, the directional derivative admits the representation $h'(x;\, d) = \langle \nabla f(x), d \rangle - g'(x;\, d)$ (or $f'(x;\, d) - g'(x;\, d)$ in the DC case), and continuity in $d$ follows from convexity \citep[Theorems~23.1 and~10.1]{Rockafellar1997}.
% \end{remark}


% \subsection{Notation}

% Let $L := \aff(C - x^0)$ denote the linear subspace parallel to $\aff(C)$. This is independent of the choice of $x^0 \in C$. For $x \in C$ and $v \in \R^n$, let
% \[
%   \tau(x;\, v) := \sup\!\bigl\{t \geq 0 : x + tv \in C\bigr\}
% \]
% denote the maximal feasible step length at $x$ in direction $v$. Fix parameters $\gamma > 0$ and $r \in (0, \infty]$ throughout.

% \begin{definition}[Cone of feasible directions]\label{def:feasible-cone}
% Let $C \subseteq \R^n$ be convex and $\bar{x} \in C$. The \emph{cone of feasible directions} is $D_C(\bar{x}) := \{w \mid \exists\, \lambda > 0,\;\bar{x} + \lambda w \in C\}$.
% \end{definition}

% The cone $D_C(\bar{x})$ is convex. A point $x^* \in C$ is \emph{d-stationary} if $h'(x^*;\, d) \geq 0$ for all $d \in D_C(x^*)$.

\section{The Random Exploration Procedure}\label{sec:meta}
The Random Exploration Procedure (REP) augments a base optimization method with a simple randomized exploration step. 
At the current iterate, REP samples a direction in the affine hull and a candidate step size, accepting the move only if it yields a feasible sufficient decrease in a regularized objective; it does not rely on gradients or problem structure and uses only function value comparisons.
This regularization guarantees sufficient-decrease-type behavior.

% The  Random Exploration Procedure (REP) couples an arbitrary inner oracle $\mathcal{O}$ with a randomized exploration step. The oracle may represent any deterministic or randomized subroutine (a spanning-directions method, a proximal step, a gradient-based update, or even an adversarial one), and is required only to produce a feasible next iterate.

\begin{algorithm}
\caption{Random Exploration Procedure (REP)}
\label{alg:ERP}
\begin{algorithmic}[1]
\Require $x \in C$; $\gamma > 0$, $r \in (0,\infty)$.
  \State Sample $v_{ex} \sim \mathcal{D}_v$ with  $\mathrm{supp}(\mathcal{D}_v) = S^{n-1} \cap L$
  \State Sample $\hat{t} \sim \mathcal{D}_t$ with  $\mathrm{supp}(\mathcal{D}_t) = [0,r]$
  \State $t_{ex} \leftarrow 
    \begin{cases}
    \hat{t} & \text{if } x + \hat{t} v_{ex} \in C \text{ and } h(x + \hat{t} v_{ex}) + \frac{\gamma}{2}(\hat{t})^2 < h(x), \\
    0 & \text{otherwise.}
    \end{cases}$
    \State \textbf{Return}  $y \leftarrow x + t_{ex} v_{ex}$
  % \State $y \leftarrow \arg\min \{h(u) : u\in \{x, x + t_{ex} v_{ex}\} \}$
\end{algorithmic}
\end{algorithm}

\begin{remark}[Sampling distributions]
The direction distribution $\mathcal{D}_v$ is supported on $S^{n-1} \cap L$ and assigns strictly positive probability to every measurable subset of $S^{n-1} \cap L$ with strictly positive surface measure. The step-size distribution $\mathcal{D}_t$ is supported on $[0,r]$ and assigns positive probability to every interval $[a,b] \subseteq (0,r]$.
We assume $(v_{\mathrm{ex}}^k,\hat t^k)$ are i.i.d.\ and independent of the past.
\end{remark}

\Cref{alg:random} augments a feasible oracle with a randomized exploration step: the oracle proposes a candidate, while exploration samples a direction and step size to generate a competing feasible point accepted by value comparison. Almost surely, every neighborhood of relative-interior feasible direction at any accumulation point is sampled infinitely often, yielding d-stationarity via Borel--Cantelli.

% The implementation in \Cref{alg:random} combines an oracle step with randomized exploration. 
% The oracle produces a feasible candidate and governs practical performance, while the exploration step forms another feasible candidate by sampling a direction and a step size, followed by a function value comparison. 
% Almost surely, every relative-interior feasible direction at any accumulation point is explored infinitely often along the outer iterations, which yields d-stationarity via a Borel--Cantelli argument.

\begin{algorithm}
\caption{Random exploration procedure implementation}
\label{alg:random}
\begin{algorithmic}[1]
\Require $x^0 \in C$; oracle $\mathcal{O}_C$ over $C$; $\gamma > 0$, $r \in (0,\infty)$.
\For{$k = 0, 1, 2, \ldots$}
  \State Set $y^{k+1} \leftarrow \mathrm{REP}(x^k)$ and  $z^{k+1} \leftarrow \mathcal{O}_C(x^k)$
  \State $x^{k+1} \leftarrow \arg\min_{u\,\in\,\{y^{k+1},\,z^{k+1}\}} h(u)$
\EndFor
\end{algorithmic}
\end{algorithm}

\begin{remark}[Implementation cost]
\Cref{alg:random} requires one feasibility check %($x^k + \hat{t}^k v^k_{ex} \in C$) 
and at most one function evaluation, %($h(x^k + \hat{t}^k v^k_{ex})$), 
which can be skipped when the candidate is infeasible, per exploration step. 
% This is strictly cheaper than performing an exact one-dimensional line search along the sampled direction.
\end{remark}

% \begin{remark}[No need to compute $\tau$]
% Algorithm~\ref{alg:random} does not require computing the maximal feasible step length $\tau(x^k;\, v^k_{ex})$. When the sampled candidate $x^k + \hat{t}^k v^k_{ex}$ is infeasible, the if-branch fails at the feasibility check and $t^k_{ex} = 0$ is selected automatically, exactly what would happen if $v^k_{ex}$ were not a feasible direction. This is in contrast to line-search implementations, where $\tau$ must be computed to define the feasible search interval for the line search.
% \end{remark}

\section{Convergence Analysis}\label{sec:convergence}

This section establishes that the exploration step enforces asymptotic d-stationarity, independently of the oracle, via a probabilistic covering of sampled directions and a local geometric property of the feasible set; some proofs are deferred to \Cref{sec:appendix Proofs of}.

% This section establishes that the exploration step enforces asymptotic d-stationarity, independently of the underlying oracle, by combining a probabilistic covering property of the sampled directions with a local geometric property of the feasible set; some proofs are deferred to \Cref{sec:appendix Proofs of}.

\begin{theorem}[Asymptotic d-stationarity]\label{thm:random}
Let  \Cref{ass:A1} and \Cref{ass:A2} hold, and let $\{x^k\}$ be generated by Algorithm~\ref{alg:random} with any oracle $\mathcal{O}_C$ producing feasible iterates. Then:
\begin{enumerate}
\item[(i)] $\sum_{k=0}^{\infty}(t^k_{ex})^2 < \infty$, and in particular $t^k_{ex} \to 0$.
\item[(ii)] Almost surely, every accumulation point of $\{x^k\}$ is d-stationary.
\end{enumerate}
\end{theorem}

% \begin{theorem}[Asymptotic d-stationarity]\label{thm:random}
% Let Assumptions~\ref{ass:A1}--\ref{ass:A3} hold, and let $\{x^k\}$ be generated by Algorithm~\ref{alg:random} with any oracle $\mathcal{O}_C$ producing feasible iterates. Then:
% \begin{enumerate}
% \item[(i)] $\sum_{k=0}^{\infty}(t^k_{ex})^2 < \infty$, and in particular $t^k_{ex} \to 0$.
% \item[(ii)] If $\{x^k\}$ has an accumulation point, then almost surely every accumulation point $x^*$ of $\{x^k\}$ is d-stationary.
% \end{enumerate}
% \end{theorem}

\begin{remark}[Existence of accumulation points]
A sufficient condition for the existence of an accumulation point is boundedness of the level set, $\{x \in C : h(x) \le h(x^0)\}.$
% \[
% \{x \in C : h(x) \le h(x^0)\}.
% \]
Indeed, greedy selection ensures $h(x^{k+1}) \le h(x^k)$ for all $k$, hence $h(x^k) \le h(x^0)$ and $\{x^k\}$ remains in this set. If it is bounded, then $\{x^k\}$ is bounded in $\R^n$ and admits an accumulation point by Bolzano--Weierstrass.
\end{remark}

% \begin{remark}[Existence of accumulation points]
% A sufficient condition for existence of an accumulation point is boundedness of the feasible level set
% \begin{equation*}
% \{x \in C : h(x) \le h(x^0)\}.
% \end{equation*}
% Indeed, greedy selection yields $h(x^{k+1}) \leq h(x^k)$ for all $k$, so by induction $h(x^k) \leq h(x^0)$ and therefore $x^k \in \{x \in C : h(x) \le h(x^0)\}$ for every $k$. Hence $\{x^k\}$ is a bounded sequence in $\R^n$, and by Bolzano--Weierstrass it admits an accumulation point.
% \end{remark}

% \subsection{Auxiliary Lemmas}
The proof of \Cref{thm:random} rests on probabilistic and geometric ingredients.
Let $\mathcal{Q}\subseteq S^{n-1}\cap L$ be any fixed countable dense set.
% For a countable dense set $\mathcal{Q} \subseteq S^{n-1} \cap L$.
Lemma \ref{lem:BC-directions} records a joint density property for the sampled exploration direction, its proximity to a target direction, and the random step size, simultaneously for all directions from $\mathcal{Q}$. Lemma~\ref{lem:step-bound} exposes a local uniform feasibility radius: for every relative-interior direction $v$, nearby directions remain in the relative interior and share a common feasible step size.% $\bar{\varepsilon}$.

\begin{lemma}[Joint density of exploration direction, proximity, and step size]\label{lem:BC-directions}
There exists an event $\Omega_{BC}$ with $\mathbb{P}(\Omega_{BC})=1$ such that, on $\Omega_{BC}$, for every $0 < a < b \leq r$, every $\rho > 0$, every accumulation point $x^*$ of $\{x^k\}$, and every $v \in \mathcal{Q}$, there exists a subsequence $J \subseteq \N$ such that, as $j \to \infty$ along $J$, it holds that: $x^j \to x^*,\ v^j_{ex} \to v,\ \norm{v^j_{ex} - v} < \rho,\ \hat{t}^j \in [a,\, b].$
% \[
%   x^j \to x^*,\qquad v^j_{ex} \to v,\qquad \norm{v^j_{ex} - v} < \rho,\qquad \hat{t}^j \in [a,\, b].
% \]
\end{lemma}

\begin{proof}
We assume throughout that $\dim L \ge 1$; otherwise $C$ is a singleton and the result is trivial.

Let $\mathcal{F}_{k-1}$ denote the $\sigma$-algebra generated by $\{x^0, v^0_{ex}, \hat{t}^0, \ldots, x^{k-1}, v^{k-1}_{ex}, \hat{t}^{k-1}, x^k\}$, so that $x^k$ is $\mathcal{F}_{k-1}$-measurable while both $v^k_{ex}$ and $\hat{t}^k$ are independent of $\mathcal{F}_{k-1}$ (and of each other).

Fix a point $y \in C$, a radius $\delta > 0$, a spherical cap
$V := B[w, \varepsilon'] \cap S^{n-1} \cap L,$
for some $w \in S^{n-1} \cap L$ and $\varepsilon' > 0$, and an interval $[a', b'] \subseteq (0, r]$ with $a' < b'$. 
Set
\[
p := \mathbb{P}(v^k_{ex} \in V) > 0,\qquad
q := \mathbb{P}(\hat t^k \in [a', b']) > 0,
\]
where positivity follows from the full-support assumptions.

Define
\[
A_k := \{v^k_{ex} \in V\} \cap \{\hat{t}^k \in [a', b']\} \cap \{\norm{x^k - y} < \delta\}.
\]
Independence yields $\mathbb{P}(A_k \mid \mathcal{F}_{k-1})
= \mathbf{1}_{\{\norm{x^k - y} < \delta\}} \cdot p \cdot q.$
% \[
% \mathbb{P}(A_k \mid \mathcal{F}_{k-1})
% = \mathbf{1}_{\{\norm{x^k - y} < \delta\}} \cdot p \cdot q.
% \]

By the conditional Borel--Cantelli--L\'evy lemma \citep[Ch.~VII, \S5, Cor.~2]{Shiryaev2019},
\begin{equation}\label{eq:BCL-step}
  \Bigl\{\sum_{k} \mathbb{P}(A_k \mid \mathcal{F}_{k-1}) = \infty\Bigr\}
  =
  \{A_k\ \text{occurs infinitely often}\}
  \quad\text{a.s.}
\end{equation}
In particular, whenever $x^k$ visits $B(y,\delta)$ infinitely often, $A_k$ occurs infinitely often almost surely.

Let $\{y_n\}$ be a countable dense subset of $C$, let $\mathcal{Q}=\{q_i\}_{i\in\N}$, and let $(\delta_m, \varepsilon'_l, a'_s, b'_s)$ range over positive rationals with $a'_s < b'_s \leq r$.
For each tuple $u=(n,i,m,l,s)$, define the event
\[
E_u := \Bigl\{\|x^k-y_n\|<\delta_m\ \text{i.o.}\ \Longrightarrow\
\bigl(v^k_{ex}\in B[q_i,\varepsilon'_l]\cap S^{n-1}\cap L,\ \hat t^k\in[a'_s,b'_s]\bigr)\ \text{i.o.}\Bigr\}.
\]
By \eqref{eq:BCL-step}, $\mathbb{P}(E_u)=1$ for every $u$: on paths where $\|x^k-y_n\|<\delta_m$ occurs infinitely often, the enriched event occurs infinitely often almost surely; otherwise the implication is vacuous.
Set $\Omega_{BC}:=\bigcap_{u} E_u.$
% \[
% \Omega_{BC}:=\bigcap_{u} E_u.
% \]
Since the index set $\{u\}$ is countable and each $E_u$ has probability $1$, we have $\mathbb{P}(\Omega_{BC})=1$. 
% \[
% \mathbb{P}(\Omega_{BC})=1.
% \]
On $\Omega_{BC}$, the following grid property holds simultaneously:

For every tuple $(y_n, \delta_m, q_i, \varepsilon'_l, a'_s, b'_s)$, if $\norm{x^k - y_n} < \delta_m$ for infinitely many $k$, then there are infinitely many such $k$ with additionally
\[
v^k_{ex} \in B[q_i, \varepsilon'_l] \cap S^{n-1} \cap L,
\qquad
\hat{t}^k \in [a'_s, b'_s].
\]

Fix $\omega \in \Omega_{BC}$. Let $x^*$ be any accumulation point of $\{x^k(\omega)\}$ and fix $v \in \mathcal{Q}$. For this fixed $\omega$, write $x^k:=x^k(\omega)$ and similarly for sampled variables. Fix arbitrary $\eta \in (0,\rho)$. By density, choose rationals $y_n, \delta_m, \varepsilon'_l, a'_s, b'_s$ such that
$
\|y_n - x^*\| < \eta/4,
\delta_m \in (\eta/2, 3\eta/4),
\varepsilon'_l < \eta/4,
$
and
$
a < a'_s < b'_s < b.
$
Choose $i$ such that $q_i=v$.
Since $x^*$ is an accumulation point, $\|x^k - x^*\| < \eta/4$ for infinitely many $k$. For such indices,
\[
\|x^k - y_n\|
\le \|x^k - x^*\| + \|x^* - y_n\|
< \eta/2 < \delta_m.
\]
Hence, by the grid property, there are infinitely many indices $k$ such that: $\|x^k - y_n\| < \delta_m, v^k_{ex} \in B[q_i, \varepsilon'_l],$ for $\hat{t}^k \in [a'_s, b'_s].$
% \[
% \|x^k - y_n\| < \delta_m,\quad
% v^k_{ex} \in B[w_i, \varepsilon'_l],\quad
% \hat{t}^k \in [a'_s, b'_s].
% \]
For these indices,
\[
\|x^k - x^*\|
\le \|x^k - y_n\| + \|y_n - x^*\|
< \delta_m + \eta/4
< \eta,
\]
and $\|v^k_{ex} - v\|
< \varepsilon'_l
< \eta < \rho,$
% \[
% \|v^k_{ex} - v\|
% < \varepsilon'_l
% < \eta < \rho,
% \]
while $\hat{t}^k \in [a'_s, b'_s] \subseteq (a,b)$.

Let $\eta_q \downarrow 0$ with $\eta_1 < \rho$, and for each $q$ let $K_q$ be the infinite set of indices satisfying the above properties with $\eta=\eta_q$. Choose inductively $j_q \in K_q$ with $j_q > j_{q-1}$. Then, since $\|x^{j_q} - x^*\| < \eta_q$ and $\eta_q \to 0$, we have $x^{j_q} \to x^*$. Similarly, $\|v^{j_q}_{ex} - v\| < \eta_q$ implies $v^{j_q}_{ex} \to v$, while $\|v^{j_q}_{ex} - v\| < \rho$ and $\hat{t}^{j_q} \in [a,b]$ for all $q$. This completes the proof.
\end{proof}


\begin{lemma}[Uniform feasibility radius]\label{lem:step-bound}
Let $x^* \in C$ and $v \in \ri(D_C(x^*)) \cap L$. There exist $\rho > 0$, $\bar{\varepsilon} > 0$, and $\delta > 0$ such that: (i)  $B(v,\, \rho) \cap L \subseteq \ri(D_C(x^*))$; (ii) $\tau(x;\, w) \geq \bar{\varepsilon}$ for every $x \in C$ with $\norm{x - x^*} < \delta$ and every $w \in B(v,\, \rho) \cap L$.
% \begin{enumerate}
% \item[(i)] $B(v,\, \rho) \cap L \subseteq \ri(D_C(x^*))$;
% \item[(ii)] $\tau(x;\, w) \geq \bar{\varepsilon}$ for every $x \in C$ with $\norm{x - x^*} < \delta$ and every $w \in B(v,\, \rho) \cap L$.
% \end{enumerate}
\end{lemma}

% \begin{proof}
% Let $k := \dim(L)$. Since $v \in \ri(D_C(x^*))$, there exists $r_0 > 0$ such that
% \[
% B(v, r_0) \cap L \subseteq \ri(D_C(x^*)).
% \]
% Pick $k+1$ directions $d_0, \ldots, d_k \in B(v, r_0) \cap L$ whose convex hull is a $k$-dimensional simplex in $L$ containing $v$ in its relative interior. Then
% \[
% \conv\{d_0, \ldots, d_k\} \subseteq \ri(D_C(x^*)).
% \]
% Since each $d_j \in D_C(x^*)$, choose $\varepsilon_j > 0$ with $x^* + \varepsilon_j d_j \in C$, and set $\bar{\varepsilon} := \min_j \varepsilon_j > 0$. By convexity of $C$, $p_j := x^* + \bar{\varepsilon}\, d_j \in C$ for each $j$, and for any convex combination $d = \sum_j \lambda_j d_j$,
% \[
%   x^* + \bar{\varepsilon}\, d = \sum_{j} \lambda_j\, p_j \in C.
% \]
% Since $v$ lies in the relative interior of the simplex, there exists $\rho_0 > 0$ with
% \[
% B(v, \rho_0) \cap L \subseteq \conv\{d_0, \ldots, d_k\},
% \]
% and hence $x^* + \bar{\varepsilon}\, w \in C$ for every $w \in B(v, \rho_0) \cap L$. The map $w \mapsto x^* + \bar{\varepsilon}\, w$ sends $B(v,\rho_0)\cap L$ onto
% \[
% B(x^* + \bar{\varepsilon}\, v,\; \bar{\varepsilon}\rho_0)\cap \aff(C),
% \]
% and therefore
% \[
%   B\!\left(x^* + \bar{\varepsilon}\, v,\; \bar{\varepsilon}\rho_0\right) \cap \aff(C) \subseteq C.
% \]

% Set $\rho := \rho_0 / 2$ and $\delta := \bar{\varepsilon}\rho_0 / 2$. Then (i) holds because
% \[
% B(v,\, \rho) \cap L \subseteq B(v,\, \rho_0) \cap L \subseteq \conv\{d_0, \ldots, d_k\} \subseteq \ri(D_C(x^*)).
% \]
% For (ii), fix $x \in C$ with $\norm{x - x^*} < \delta$ and $w \in B(v,\, \rho) \cap L$. Then
% \[
%   \norm{(x + \bar{\varepsilon}\, w) - (x^* + \bar{\varepsilon}\, v)} \leq \norm{x - x^*} + \bar{\varepsilon}\norm{w - v} < \delta + \bar{\varepsilon}\rho = \bar{\varepsilon}\rho_0,
% \]
% and since $x \in C \subseteq \aff(C)$ and $w \in L$, we have $x + \bar{\varepsilon}\, w \in \aff(C)$; hence
% \[
% x + \bar{\varepsilon}\, w \in B(x^* + \bar{\varepsilon}\, v,\, \bar{\varepsilon}\rho_0) \cap \aff(C) \subseteq C.
% \]
% Convexity then gives
% \[
%   x + t w = (1 - t/\bar{\varepsilon})\, x + (t/\bar{\varepsilon})\,(x + \bar{\varepsilon}\, w) \in C \qquad \text{for every } t \in [0,\, \bar{\varepsilon}],
% \]
% so $\tau(x;\, w) \geq \bar{\varepsilon}$.
% \end{proof}

\begin{lemma}[Continuity and extension]\label{lem:continuity}
The directional derivative map $d \mapsto h'(x^*;\, d)$ is continuous and positively homogeneous of degree one. Moreover, if $h'(x^*;\, v) \geq 0$ for all $v \in \ri(D_C(x^*))$, then $h'(x^*;\, d) \geq 0$ for all $d \in \cl(D_C(x^*))$.
\end{lemma}

% \begin{proof}
% Let $K := D_C(x^*)$, a convex cone. By Assumption~A3, the map $d \mapsto h'(x^*;\, d)$ is continuous.
% We next show positive homogeneity. For any $\alpha > 0$,
% \[
%   h'(x^*;\, \alpha d)
%   = \lim_{t\downarrow 0}\frac{h(x^* + t\alpha d)-h(x^*)}{t}
%   = \alpha \lim_{t\downarrow 0}\frac{h(x^* + t d)-h(x^*)}{t}
%   = \alpha h'(x^*;\, d).
% \]

% For the extension statement, recall that for any convex set $K$,
% \[
% \cl(\ri(K)) = \cl(K),
% \]
% where the relative interior is taken with respect to $\aff(K)=L$. Let $d \in \cl(K)$. Then there exists a sequence $v_m \in \ri(K)$ with $v_m \to d$. By continuity,
% \[
%   h'(x^*;\, d) = \lim_{m\to\infty} h'(x^*;\, v_m) \geq 0,
% \]
% since $h'(x^*;\, v_m) \geq 0$ for all $m$.
% \end{proof}


The probabilistic and geometric infrastructure is now ready for the proof of \Cref{thm:random}.

% \subsection{Main Theorem}

\begin{proof}[Proof of \Cref{thm:random}]
\textit{(i) Step-size decay.} Since $t = 0$ is always available in the two-point comparison and $t^k_{ex}$ minimizes over $\{0, \hat{t}^k\}$:
  $h\!\left(x^k + t^k_{ex}\,v^k_{ex}\right) + \tfrac{\gamma}{2}(t^k_{ex})^2 \leq h(x^k).$
  
By greedy selection, $h(x^{k+1}) \leq h(y^{k+1})$, so $h(x^k) - h(x^{k+1}) \geq \frac{\gamma}{2}(t^k_{ex})^2$. Summing gives $\sum_k (t^k_{ex})^2 \leq \frac{2}{\gamma}(h(x^0) - \inf_C h) < \infty$, hence $t^k_{ex} \to 0$.

\textit{(ii) D-stationarity.} Let $\Omega_{BC}$ be the event from Lemma~\ref{lem:BC-directions}, which guarantees simultaneous subsequences for all directions in $\mathcal{Q}$. Fix $\omega \in \Omega_{BC}$ and work pathwise. Let $x^*$ be an accumulation point of $\{x^k(\omega)\}$, and enumerate $\mathcal{Q}\cap \ri(D_C(x^*))$ as $\{v_i\}_{i=0}^{\infty}$. Since $\ri(D_C(x^*))\cap S^{n-1}$ is relatively open in $S^{n-1}\cap L$ and $\mathcal{Q}$ is dense in $S^{n-1}\cap L$, the set $\{v_i\}$ is dense in $\ri(D_C(x^*))\cap S^{n-1}$.

Fix $i \in \N$ and suppose for contradiction $h'(x^*;\, v_i) = -\varepsilon < 0$ for some $\varepsilon > 0$.



By Lemma~\ref{lem:step-bound} applied with $v = v_i$, there exist $\rho_i > 0$, $\bar{\varepsilon}_i > 0$, and $\delta_i > 0$ such that $B(v_i,\, \rho_i) \cap L \subseteq \ri(D_C(x^*))$, and $\tau(x;\, w) \geq \bar{\varepsilon}_i$ for every $x \in C$ with $\norm{x - x^*} < \delta_i$ and every $w \in B(v_i,\, \rho_i) \cap L$.



By local Lipschitz continuity of $h$ at $x^*$, there exist $R_h>0$ and $L_h>0$ such that $h$ is $L_h$-Lipschitz on $B(x^*,R_h)$.

Next choose the comparison scale. By definition of the directional derivative with $h'(x^*;\, v_i) = -\varepsilon$, there exists $\tau_0 > 0$ such that $h(x^* + t v_i) - h(x^*) \leq -\tfrac{\varepsilon}{2}\,t$ for every $t \in (0,\, \tau_0]$. Set
\[
  t^* := \min\!\left\{\bar{\varepsilon}_i,\;\tau_0,\;\frac{\varepsilon}{4\gamma},\;r,\;\frac{R_h}{4}\right\} > 0.
\]



Now lock direction and step-size sampling simultaneously. Since $\omega \in \Omega_{BC}$ and the parameters are now fixed, Lemma~\ref{lem:BC-directions} gives an infinite subsequence $\mathcal{J}_i \subseteq \N$ along which $x^j \to x^*,  v^j_{ex} \to v_i,  \norm{v^j_{ex} - v_i} < \rho_i,  \hat{t}^j \in [t^*/2,\, t^*].$

For $j \in \mathcal{J}_i$ sufficiently large, $\norm{x^j - x^*} < \delta_i$; the feasibility bound above then gives $\tau(x^j;\, v^j_{ex}) \geq \bar{\varepsilon}_i \geq t^*$, so $x^j + t v^j_{ex} \in C$ for every $t \in [0,\, t^*]$.



The key point is that the same strict decrease estimate holds uniformly for every $t \in [t^*/2,\, t^*]$. 
Since $x^j \to x^*$ along $\mathcal{J}_i$, for all sufficiently large $j \in \mathcal{J}_i$ we have $\norm{x^j-x^*}<R_h/4$. For such $j$ and any $t\in[0,t^*]$, $\norm{x^*+t v_i-x^*}=t\le t^*\le R_h/4<R_h,$
% \[
% \norm{x^*+t v_i-x^*}=t\le t^*\le R_h/4<R_h,
% \]
and, using $\norm{v_{ex}^j}=1$,
\[
\norm{x^j+t v_{ex}^j-x^*}\le \norm{x^j-x^*}+t\,\norm{v_{ex}^j}
<R_h/4+t^*\le R_h/2<R_h.
\]
Hence both points $x^j+t v_{ex}^j$ and $x^*+t v_i$ lie in $B(x^*,R_h)$, so the $L_h$-Lipschitz bound is valid. Therefore, for any $t \in [t^*/2,\, t^*]$ and all sufficiently large $j \in \mathcal{J}_i$, decompose
\begin{align*}
  h(x^j + t v^j_{ex}) - h(x^j) + \tfrac{\gamma}{2}t^2
  &= \bigl[h(x^j + t v^j_{ex}) - h(x^* + t v_i)\bigr]
   + \bigl[h(x^* + t v_i) - h(x^*)\bigr] \\
  &\quad + \bigl[h(x^*) - h(x^j)\bigr] + \tfrac{\gamma}{2}t^2 \\
  &\leq L_h\bigl(\norm{x^j - x^*} + t\,\norm{v^j_{ex} - v_i}\bigr)
     - \tfrac{\varepsilon}{2}t
     + L_h\norm{x^j - x^*}
     + \tfrac{\gamma}{2}t^2.
\end{align*}
Since $t \in [t^*/2,\, t^*]$, $\tfrac{\varepsilon}{2}t \geq \tfrac{\varepsilon t^*}{4}$; and $t^* \leq \varepsilon/(4\gamma)$ gives $\tfrac{\gamma}{2}t^2 \leq \tfrac{\gamma}{2}(t^*)^2 \leq \tfrac{\varepsilon t^*}{8}$. 
Since $x^j \to x^*$ and $v^j_{ex} \to v_i$ along $\mathcal{J}_i$, and since $[t^*/2,\, t^*]$ is compact and $x^j \to x^*$, $v^j_{ex} \to v_i$, we may pass to a tail subsequence (still denoted $\mathcal{J}_i$) on which, uniformly in $t \in [t^*/2,\, t^*]$,
\[
  L_h\bigl(\norm{x^j - x^*} + t\,\norm{v^j_{ex} - v_i}\bigr) + L_h\norm{x^j - x^*}
  \;\leq\; 2 L_h\norm{x^j - x^*} + L_h\, t^*\norm{v^j_{ex} - v_i}
  \;\leq\; \tfrac{\varepsilon t^*}{16}.
\]
Therefore, for every $t \in [t^*/2,\, t^*]$ we have that
\[
  h(x^j + t v^j_{ex}) - h(x^j) + \tfrac{\gamma}{2}t^2
  \;\leq\; \tfrac{\varepsilon t^*}{16} - \tfrac{\varepsilon t^*}{4} + \tfrac{\varepsilon t^*}{8}
  \;=\; -\tfrac{\varepsilon t^*}{16} \;<\; 0.
\]
Applied at $t = \hat{t}^j \in [t^*/2,\, t^*]$, together with the feasibility established above, this strict negativity forces the two-point comparison to select $t^j_{ex} = \hat{t}^j \geq t^*/2$ for every $j \in \mathcal{J}_i$.


This yields the contradiction: for infinitely many $j \in \mathcal{J}_i$, $t^j_{ex} \geq t^*/2$, so $\sum_k (t^k_{ex})^2 \geq \sum_{j \in \mathcal{J}_i}(t^*/2)^2 = +\infty$, contradicting part~(i). Hence $h'(x^*;\, v_i) \geq 0$.


Since $\{v_i\}$ is dense in $\ri(D_C(x^*)) \cap S^{n-1}$, Lemma~\ref{lem:continuity} (continuity and positive homogeneity of $h'(x^*;\, \cdot)$, together with the relative-interior-to-closure extension) yields $h'(x^*;\, d) \geq 0$ for every $d \in \cl(D_C(x^*))$. Hence $x^*$ is d-stationary. As $x^*$ was an arbitrary accumulation point on this fixed $\omega \in \Omega_{BC}$, every accumulation point is d-stationary on $\Omega_{BC}$, proving the almost-sure claim.
\end{proof}

The preceding argument requires only that the oracle outputs feasible points $z^{k+1} \in C$; otherwise, $\mathcal{O}_C$ may be arbitrary, while d-stationarity is enforced solely by the exploration step.



\section{Oracle Instantiations}\label{sec:oracles}
% This section illustrates the oracle--augmentation separation through three representative instantiations: DCA under a strongly convex DC model (Section~\ref{sec:dc}), prox-linear-type oracles for nonsmooth composite objectives $h=c\circ F$ on $\R^n$ (Section~\ref{subsec:pl-oracle}), and gradient sampling via INGD (Section~\ref{subsec:gs-oracle}).
% In each case, oracle-specific structure supplies quantitative performance (a stationarity-gap rate or an INGD complexity bound), while Theorem~\ref{thm:random} supplies almost-sure asymptotic d-stationarity of accumulation points through the exploration step alone.

We illustrate oracle-driven rates versus exploration-driven optimality on DCA, prox-linear methods, and gradient sampling (INGD): in each case the oracle supplies a stationarity measure and rate, while Theorem~\ref{thm:random} upgrades accumulation points to d-stationarity almost surely.

% \Cref{alg:random} can utilize various oracles.
% We now turn to provide three prominent examples.
% The first is the classical and prototypical framework of DCA, achieving (DC) critical points and not d-stationary points.

\subsection{DCA for constrained DC programming}\label{sec:dc}
Consider the constrained difference-of-convex problem
\begin{equation}\label{eq:dca}
    \min_{x \in C} h(x) := f_1(x) - f_2(x),
\end{equation}
where $C \subseteq \R^n$ is closed and convex,  $f_1$ is $\mu_1$-strongly convex and $f_2$ is $\mu_2$-strongly convex with $\mu_1 + \mu_2 > 0$, 
and that $h$ is bounded below.

% The strong convexity assumptions on $f_1$ and $f_2$ imposed below are strictly stronger than Assumptions~A1--A3 and are used only to obtain the quantitative rate in Theorem~\ref{thm:dca-rate}; the argument needs no smoothness of $f_1$, only the DCA subproblem optimality relation and the strong convexity inequality at subgradients. Almost-sure d-stationarity of accumulation points as in Theorem~\ref{thm:random} remains a consequence of Assumptions~A1--A3 and the exploration step, independently of this oracle.

% Constrained difference-of-convex (CD) programming problem is given by:
% \begin{equation}\label{eq:dca}
%   \min_{x \in C} h(x) := f_1(x) - f_2(x),
% \end{equation}
% where $C \subseteq \R^n$ is nonempty, closed, and convex; $f_1$ is proper, closed, and $\mu_1$-strongly convex (and possibly non-differentiable); $f_2$ is $\mu_2$-strongly convex, with $\mu_2 \geq 0$; $\mu_1 + \mu_2 > 0$; $h^* := \inf_{x \in C} h(x) > -\infty$. 
% The constraint set $C$ is incorporated via the indicator function $\delta_C(x) = 0$ if $x \in C$ and $+\infty$ otherwise, so the effective convex part in the DC decomposition is $\tilde{f}_1 := f_1 + \delta_C$. 
% For $x \in C$, $\partial \delta_C(x) = N_C(x)$, the normal cone to $C$ at $x$, and therefore
% $\partial \tilde{f}_1(x) = \partial f_1(x) + N_C(x).$

% The constraint set $C$ is incorporated via the indicator function $\delta_C(x) = 0$ if $x \in C$ and $+\infty$ otherwise, so the effective convex part in the DC decomposition is $\tilde{f}_1 := f_1 + \delta_C$. 
% For $x \in C$, $\partial\delta_C(x) = N_C(x)$, the normal cone to $C$ at $x$, so $\partial\tilde{f}_1(x) = \nabla f_1(x) + N_C(x)$.

% \begin{remark}[Relation to the general setting]
% The assumptions of this section are strictly stronger than Assumptions~A1--A3: strong convexity of $f_1$ and $f_2$ is not required elsewhere in the paper. They are imposed here to obtain a quantitative rate. Notably, $L_1$-smoothness (and indeed any differentiability) of $f_1$ is not required: the proof uses only the subdifferential optimality condition~\eqref{eq:OC} and the strong convexity definition applied at a subgradient. The d-stationarity guarantee of Theorem~\ref{thm:random} holds under the weaker Assumptions~A1--A3 alone.
% \end{remark}

% In DC programming, the baseline optimality condition is the so-called \textit{criticality} condition.
% A point $\bar{x} \in C$ is \emph{critical} for the constrained DC problem if $\partial\tilde{f}_1(\bar{x}) \cap \partial f_2(\bar{x}) \neq \emptyset$, i.e., there exist $g_1\in \partial f_1(\bar x)$, $g_2\in \partial f_2(\bar x)$, and $n\in N_C(\bar x)$ such that $g_1+n=g_2$.
% % there exist $g_2 \in \partial f_2(\bar{x})$ and $n \in N_C(\bar{x})$ such that $\nabla f_1(\bar{x}) + n = g_2$.

% \subsection{The DCA Oracle}\label{subsec:dca-oracle}

Given $x^k \in C$, the DCA oracle produces
\[
z^{k+1} \in \arg\min_{x \in C} \; f_1(x) - \langle g_2^k, x \rangle, 
\quad g_2^k \in \partial f_2(x^k).
\]
This step characterizes criticality through the fixed-point condition $x^k = z^{k+1}$.


% The Difference-of-Convex Algorithm (DCA) oracle produces, given the current iterate $x^k \in C$, a new feasible point
% \[
%   z^{k+1} = \arg\min_{x \in C}\; f_1(x) - \inner{g^k_2}{x}, \qquad g^k_2 \in \partial f_2(x^k).
% \]
% The first-order optimality condition for this strongly convex subproblem is:
% \begin{equation}
%   \exists\; s^{k+1} \in \partial f_1(z^{k+1}),\quad n^{k+1} \in N_C(z^{k+1}) \quad\text{such that}\quad s^{k+1} + n^{k+1} = g^k_2. \tag{OC}\label{eq:OC}
% \end{equation}
% In particular, $x^k - z^{k+1} = 0$ if and only if $z^{k+1} = x^k$, which by~\eqref{eq:OC} is exactly the criticality condition for $x^k$.

% \subsection{Algorithm}\label{subsec:dca-algorithm}

The augmented method applies Algorithm~\ref{alg:random} with this oracle. 
The key observation is that the rate analysis of DCA relies only on the descent inequality
$h(x^{k+1}) \leq h(z^{k+1}),$
which is preserved by the greedy selection step regardless of whether the exploration step is accepted; the proof is deferred to Appendix~\ref{sec:appendix-dca}.


% The algorithm is Algorithm~\ref{alg:random} with $\mathcal{O}(x^k) = z^{k+1}$ as described above. For convenience, we record the key consequence of the two-point comparison and greedy selection: by the argument of Section~\ref{sec:convergence},
% \begin{equation}
%   h(x^{k+1}) \leq h(y^{k+1}) \leq h(x^k) - \tfrac{\gamma}{2}(t^k_{ex})^2, \tag{G}\label{eq:G}
% \end{equation}
% and in particular $h(x^{k+1}) \leq h(z^{k+1})$ and $\{h(x^k)\}$ is non-increasing.

% \subsection{Convergence Rate Proof}\label{subsec:dca-rate}

\begin{theorem}[Rate preservation under DCA]\label{thm:dca-rate}
Under the assumptions above, after $N$ iterations of Algorithm~\ref{alg:random} with the DCA oracle, $\min_{0 \leq k \leq N-1} \norm{x^k - z^{k+1}}
\leq 
\sqrt{\frac{2\,(h(x^0) - h^*)}{(\mu_1 + \mu_2)\,N}}.$
% \[
% \min_{0 \leq k \leq N-1} \norm{x^k - z^{k+1}}
% \leq 
% \sqrt{\frac{2\,(h(x^0) - h^*)}{(\mu_1 + \mu_2)\,N}}.
% \]
\end{theorem}

% \begin{theorem}[Convergence rate]\label{thm:dca-rate}
% Let the assumptions stated above  hold. 
% After $N$ steps of Algorithm~\ref{alg:random} with the DCA oracle:
% \[
%   \min_{0 \leq k \leq N-1} \norm{x^k - z^{k+1}} \leq \sqrt{\frac{2\,(h(x^0) - h^*)}{(\mu_1 + \mu_2)\,N}}.
% \]
% \end{theorem}

% \begin{proof}
% See Appendix~\ref{sec:appendix-dca}.
% \end{proof}

Thus, the augmentation preserves the $O(1/\sqrt{N})$ rate on the natural oracle residual, while Theorem~\ref{thm:random} ensures that every accumulation point is almost surely d-stationary, which is strictly stronger than the criticality guarantee of DCA.
That is, the oracle governs the convergence rate and identifies critical points, while the exploration step 
eliminates spurious stationary points by ruling out feasible descent directions. 
These two mechanisms operate independently.

% \begin{remark}[Constrained DC via the indicator and unconstrained-rate comparison]
% Passing from unconstrained DCA to~\eqref{eq:dca} replaces $f_1$ by $\tilde{f}_1 = f_1 + \delta_C$, which introduces $n^{k+1}$ in~\eqref{eq:OC}; the normal cone inequality $\inner{n^{k+1}}{x^k - z^{k+1}} \leq 0$ removes this term from the descent lemma for $x^k \in C$. No differentiability of $f_1$ is required: $s^{k+1} \in \partial f_1(z^{k+1})$ enters as in the classical analysis. In unconstrained analyses one may track a gradient criticality gap of the form $\norm{g^k_1 - g^k_2}$; under constraints this quantity is not a faithful stationarity residual, because normal cone terms in~\eqref{eq:OC} can be large independently of how close $x^k$ is to a critical point, whereas $\norm{x^k - z^{k+1}}$ matches the fixed-point characterization of the constrained oracle step. For unconstrained DCA across several regularity regimes, \citet{Abbaszadehpeivasti2024} develop semidefinite programming performance estimation and obtain sharp constant-level guarantees in their chosen gap-type metrics; that line of work is complementary, and we do not aim to sharpen constants relative to their unconstrained benchmarks. Theorem~\ref{thm:dca-rate} records only an $O(1/\sqrt{N})$ scale bound on $\min_k \norm{x^k - z^{k+1}}$ under the strongly convex hypotheses used here, proved by strong convexity and telescoping.
% \end{remark}

% \subsection{Discussion}\label{subsec:dca-discussion}

% \begin{remark}[Oracle residual $\norm{x^k - z^{k+1}}$]\label{rem:pl-oracle-residual}\label{rem:dca-optimality-measure}
% For the DCA oracle, $\norm{x^k - z^{k+1}} = 0$ if and only if $x^k$ is critical (fixed point of the oracle subproblem). The iterate displacement $\norm{x^{k+1} - x^k}$ admits the same $O(1/\sqrt{N})$ scale but mixes oracle progress with the exploration step unless $t^k_{ex}$ is negligible or the greedy step selects $z^{k+1}$. For prox-linear-type oracles, the same quantity is the natural outer-oracle step. Appendix~\ref{sec:appendix-pl} records limit theorems (Theorems~\ref{thm:app-pl-dc}, \ref{thm:app-pl-clarke}, and~\ref{thm:app-pl-mord}) showing that, \emph{without} augmentation, decay of $\norm{x^k - z^{k+1}}$ along a convergent subsequence implies DC criticality, Clarke stationarity, or Mordukhovich stationarity of $h = c \circ F$ at the limit under templates~(DC), (C), or~(M) (Assumption~\ref{ass:app-pl-bdd}); see \citep[Theorem~8.49]{RockafellarWets2009} and Remark~\ref{rem:stationarity-comparison}. This is not in tension with Theorem~\ref{thm:random}: augmentation upgrades accumulation points to almost-sure d-stationarity, strictly stronger in general. Theorem~\ref{thm:pl-rate} gives an ergodic $O(1/\sqrt{N})$ bound in the same currency; Lemma~\ref{lem:pl-descent} links summability to the regime in which those limits apply; Remark~\ref{rem:app-pl-purpose} summarizes the appendix.
% \end{remark}

% \begin{remark}[Preservation of the rate under augmentation]\label{rem:rate-preservation}
% The proof of Theorem~\ref{thm:dca-rate} uses only two properties of Algorithm~\ref{alg:random}: (i)~$h(x^{k+1}) \leq h(z^{k+1})$, which follows from greedy selection and holds for any feasible oracle; (ii)~Lemma~\ref{lem:dca-descent}, which is a property of the DCA oracle alone. 
% The exploration step does not affect this rate proof, since the proof uses only
% $h(x^{k+1})\le h(z^{k+1})$.
% Moreover, it may strictly improve the constant $h(x^0) - h^*$: on steps where the exploration wins the greedy comparison, the running minimum of $h$ is driven below what the DCA oracle alone would achieve.

% More generally, the rate carries over to any stateless descent oracle satisfying a quadratic descent inequality
% \begin{equation}
%   h(x^k) - h(z^{k+1}) \geq \alpha\norm{x^k - z^{k+1}}^2 \tag{QD}\label{eq:QD}
% \end{equation}
% for some fixed $\alpha > 0$. The prox-linear-type oracles of Section~\ref{subsec:pl-oracle} satisfy~\eqref{eq:QD} with $\alpha = \tfrac{1}{2}(\rho - L_c L_F)$ whenever $\rho > L_c L_F$ and the model decrease condition~\eqref{eq:pl-model-decrease} holds (Lemma~\ref{lem:pl-descent}), and the argument applies verbatim.

% By contrast, the proximal bundle method, while a descent algorithm, does not satisfy the quadratic inequality in the required form: its guaranteed decrease at serious steps is expressed in terms of the cutting-plane model's predicted decrease, not a uniform quadratic in $\norm{x^k - z^{k+1}}$. Rate preservation for the bundle method as oracle therefore does not follow from the present argument.
% \end{remark}

% \begin{remark}[Oracle-side rates vs.\ exploration d-stationarity]
% Theorem~\ref{thm:dca-rate} bounds the criticality residual $\min_k \norm{x^k - z^{k+1}}$; standard closure arguments using~\eqref{eq:OC} yield critical accumulation points, a condition weaker than d-stationarity. Theorem~\ref{thm:random} upgrades every accumulation point to d-stationarity almost surely. These quantitative and geometric conclusions are logically independent (descent along the oracle vs.\ Borel--Cantelli coverage of feasible directions). For gradient sampling, INGD~\citep{Zhang2020} is analyzed in Goldstein's sense~\citep{Goldstein1977,Burke2005}: its $\tilde{O}(\varepsilon^{-3}\delta^{-1})$ complexity to reach $(\delta,\varepsilon)$-stationarity uses only that $h$ is bounded below and that accepted outer steps are monotone in $h$, so greedy selection preserves the bound (Section~\ref{subsec:gs-oracle}). That guarantee is likewise independent of Theorem~\ref{thm:random}, and Goldstein stationarity is asymptotically weaker than d-stationarity in general.
% \end{remark}

% % \section{Further Oracle Instantiations: Convergence Rates}\label{sec:oracles}

% Remark~\ref{rem:rate-preservation} observed that the rate-preservation argument of Section~\ref{sec:dc} applies verbatim to any stateless descent oracle satisfying a quadratic descent inequality~\eqref{eq:QD} for some fixed $\alpha > 0$. This section makes two further instantiations explicit. Section~\ref{subsec:pl-oracle} develops prox-linear-type oracles under a globally Lipschitz outer function $c$ (not assumed convex), where~\eqref{eq:QD} holds with $\alpha = \tfrac{1}{2}(\rho - L_c L_F)$ under the model decrease condition~\eqref{eq:pl-model-decrease}. For the gradient sampling oracle, we instead cite the existing $\tilde{O}(\varepsilon^{-3}\delta^{-1})$ complexity bound of \citet{Zhang2020} for $(\delta,\varepsilon)$-Goldstein stationarity and observe that greedy selection preserves it without modification.

\subsection{Prox-linear Oracle}\label{subsec:pl-oracle}
We next consider composite objectives of the form
\begin{equation}\label{eq:prox-linear}
    \min_{x \in \R^n} h(x) := c(F(x)),
\end{equation}
where $h$ is bounded below, $F$ is continuously differentiable with Lipschitz Jacobian and $c$ is globally Lipschitz 
and directionally differentiable. 
We assume in addition that $c$ is \emph{Hadamard directionally differentiable} on $\R^m$: for every $v, w \in \R^m$ the limit $ c'(v;\, w) := \lim_{(t,\,w') \to (0^+,\, w)} (c(v + t w') - c(v))/t$
% \[
%   c'(v;\, w) := \lim_{(t,\,w') \to (0^+,\, w)} \frac{c(v + t w') - c(v)}{t}
% \]
exists, and for each fixed $v$ the map $w \mapsto c'(v;\, w)$ is continuous on $\R^m$.
Since $F$ is continuously differentiable, standard chain rules for Hadamard directional derivatives imply that $h = c \circ F$ is locally Lipschitz (A1), Hadamard directionally differentiable at every $x \in \R^n$ with
\begin{equation}\label{eq:pl-hadamard-chain}
  h'(x;\, d) = c'\!\bigl(F(x);\, \nabla F(x)\, d\bigr),
\end{equation}
and $d \mapsto h'(x;\, d)$ exists and is continuous.
Hence, under these conditions, $h$ satisfies \Cref{ass:A1} and \Cref{ass:A2}.

% Given $x^k$, a prox-linear oracle produces $z^{k+1}$ by approximately minimizing the local model (see, e.g., \citep{Drusvyatskiy2019})
% $m_k(y) := c\big(F(x^k) + \nabla F(x^k)(y - x^k)\big) + \frac{\rho}{2}\norm{y - x^k}^2,
% $
% and enforces the model decrease condition
Given $x^k$, define the local model
$m^k(y) := c\big(F(x^k) + \nabla F(x^k)(y - x^k)\big) + \frac{\rho}{2}\norm{y - x^k}^2,$
and, because $c$ need not be convex, let the oracle return $z^{k+1}$ from an Appendix~\ref{sec:appendix-pl} first-order subproblem variant while enforcing
\begin{equation}\label{eq:pl-model-decrease}
    m^k(z^{k+1}) \leq m^k(x^k) = h(x^k).
\end{equation}

When $\rho$ is sufficiently large, this implies $h(z^{k+1}) \le h(x^k)$ and $h(x^k)-h(z^{k+1}) \ge \alpha\norm{z^{k+1}-x^k}^2$ for some $\alpha>0$.
% When $\rho$ is sufficiently large, this implies the descent inequality
% $h(z^{k+1}) \leq h(x^k),$
% and yields the quadratic decrease estimate
% $h(x^k) - h(z^{k+1}) \;\geq\; \alpha \norm{z^{k+1} - x^k}^2$
% for some $\alpha > 0$ depending on the problem constants.

Applying the same argument as in the DCA case, we obtain (proof is provided in Appendix \ref{subsec:app-pl-descent-and-rate}):

\begin{theorem}[Rate preservation for prox-linear oracles]\label{thm:pl-rate}
Suppose the oracle satisfies \eqref{eq:pl-model-decrease} and the quadratic decrease inequality above. 
Then, after $N$ iterations of Algorithm~\ref{alg:random}, $\min_{0 \leq k \leq N-1} \norm{x^k - z^{k+1}}
\leq 
\sqrt{\frac{h(x^0) - h^*}{\alpha N}}.$
% \[
% \min_{0 \leq k \leq N-1} \norm{x^k - z^{k+1}}
% \leq 
% \sqrt{\frac{h(x^0) - h^*}{\alpha N}}.
% \]
\end{theorem}
% \begin{proof}
% See Subsection~\ref{subsec:app-pl-descent-and-rate} of Appendix~\ref{sec:appendix-pl}.
% \end{proof}

As before, the exploration step does not enter the rate proof. It only modifies the iterate selection, 
while preserving the descent structure required by the oracle analysis.

The prox-linear framework includes a broad class of nonsmooth nonconvex problems; different inner solvers may yield different notions of stationarity (e.g., Clarke or Mordukhovich stationarity).
Regardless of these distinctions, Theorem~\ref{thm:random} upgrades all accumulation points of the augmented method to d-stationarity almost surely.


% --------- Old detailed version ----------\\

% The quadratic prox-linear descent step used below is standard for globally Lipschitz composite models \citep{Drusvyatskiy2019}, and we organize the discussion around three inner-termination templates compatible with a single model-decrease rule, deferring the synthesis of asymptotic limits, finite-horizon rates, and almost-sure d-stationarity to Remark~\ref{rem:pl-oracle-residual}.

% We consider
% \begin{equation}\label{eq:prox-linear}
%   \min_{x \in \R^n} h(x) := c(F(x)),
% \end{equation}
% where $F \colon \R^n \to \R^m$ is continuously differentiable with $L_F$-Lipschitz Jacobian in operator norm,
% \[
%   \norm{\nabla F(y) - \nabla F(x)}_{\mathrm{op}} \leq L_F \norm{y - x}\qquad \forall\, x,y \in \R^n,
% \]
% and $c \colon \R^m \to \R$ is globally Lipschitz with constant $L_c$.
% We fix $\rho > L_c L_F$ and write $h^* := \inf_{x \in \R^n} h(x) > -\infty$.

% For Assumptions~A1--A3 we assume in addition that $c$ is \emph{Hadamard directionally differentiable} on $\R^m$: for every $v, w \in \R^m$ the limit
% \[
%   c'(v;\, w) := \lim_{(t,\,w') \to (0^+,\, w)} \frac{c(v + t w') - c(v)}{t}
% \]
% exists, and for each fixed $v$ the map $w \mapsto c'(v;\, w)$ is continuous on $\R^m$.
% Since $F$ is continuously differentiable, standard chain rules for Hadamard directional derivatives imply that $h = c \circ F$ is locally Lipschitz (A1), Hadamard directionally differentiable at every $x \in \R^n$ with
% \begin{equation}\label{eq:pl-hadamard-chain}
%   h'(x;\, d) = c'\!\bigl(F(x);\, \nabla F(x)\, d\bigr),
% \end{equation}
% and $d \mapsto h'(x;\, d)$ is continuous (A2--A3). Indeed, A2 follows by composing the joint limit defining $c'(F(x);\,\cdot)$ with the linear map $d \mapsto \nabla F(x)\,d$; for A3, if $d^n \to d$, then $w^n := \nabla F(x)\,d^n \to \nabla F(x)\,d$, and continuity of $w \mapsto c'(F(x);\, w)$ yields $h'(x;\, d^n) = c'(F(x);\, w^n) \to c'(F(x);\, \nabla F(x)\,d) = h'(x;\, d)$. The same continuity-in-directions pattern appears for the INGD class in Section~\ref{subsec:gs-oracle}.

% Write $\ell^k(y) := F(x^k) + \nabla F(x^k)(y - x^k)$ and the prox-linear surrogate
% \[
%   m^k(y) := c(\ell^k(y)) + \tfrac{\rho}{2}\norm{y - x^k}^2.
% \]
% All variants below impose the \emph{model decrease} condition
% \begin{equation}\label{eq:pl-model-decrease}
%   m^k(z^{k+1}) \leq m^k(x^k) = h(x^k).
% \end{equation}
% This holds, for example, if the inner routine is initialized at $x^k$ and is monotone with respect to $m^k$ along its inner iterates.
% Moreover, whenever $\rho > L_c L_F$, Lipschitzness of $c$ and the second-order error bound for $\ell^k$ relative to $F$ imply the majorization $h(y) \leq m^k(y)$ for all $y \in \R^n$ \citep{Drusvyatskiy2019}.
% In particular $h(z^{k+1}) \leq m^k(z^{k+1})$, so~\eqref{eq:pl-model-decrease} yields $h(z^{k+1}) \leq h(x^k)$; together with greedy selection in Algorithm~\ref{alg:random}, $h(x^{k+1}) \leq h(z^{k+1})$ gives monotonicity $h(x^{k+1}) \leq h(x^k)$.

% Theorem~\ref{thm:pl-rate} and Lemma~\ref{lem:pl-descent} require only~\eqref{eq:pl-model-decrease} (and the Lipschitz data above); they do not depend on which first-order template is used below.
% For the asymptotic stationarity statements in Appendix~\ref{sec:appendix-pl}, we also assume throughout that the initial sublevel set
% \[
%   \{x \in \R^n : h(x) \leq h(x^0)\}
% \]
% is bounded. Under greedy selection in Algorithm~\ref{alg:random}, one then has $h(x^{k+1}) \leq h(x^k)$ for all $k$, hence $h(x^k) \leq h(x^0)$ and the iterate sequence $\{x^k\}$ remains bounded, so accumulation points exist.
% The oracle output $z^{k+1} \in \R^n$ is required to satisfy~\eqref{eq:pl-model-decrease} and one of the following first-order templates for the model $m^k$:
% \begin{enumerate}
% \item[(DC)] \textbf{DC outer.} Assume in addition that $c = f - g$ with $f,g$ convex and finite-valued on $\R^m$ (a structured subclass of the standing Lipschitz--Hadamard hypothesis on $c$). The inner step is DC critical for $m^k$ relative to the decomposition $m^k = (f \circ \ell^k + \tfrac{\rho}{2}\norm{\cdot - x^k}^2) - (g \circ \ell^k)$ (see Appendix~\ref{sec:appendix-pl}, \eqref{eq:app-pl-dc-inner}).
% \item[(C)] \textbf{Clarke outer.} Assume $\nabla F(x)$ is surjective for every $x \in \R^n$. The inner step is Clarke stationary for $y \mapsto m^k(y)$ at $y = z^{k+1}$ (Appendix~\ref{sec:appendix-pl}, \eqref{eq:app-pl-clarke-inner}).
% \item[(M)] \textbf{Mordukhovich outer.} The same surjectivity assumption as~(C). The inner step is Mordukhovich (limiting) stationary for $y \mapsto m^k(y)$ at $y = z^{k+1}$ (Appendix~\ref{sec:appendix-pl}, \eqref{eq:app-pl-mord-inner}).
% \end{enumerate}

% Classical \emph{convex-composite} prox-linear analyses \citep{Drusvyatskiy2019} typically assume a convex finite outer $c$, take $\rho \geq L_c L_F/2$ for the majorization $h \leq m^k$, and exploit strong convexity of $m^k$ together with exact or near-exact model minimization to obtain sharper stationarity and model-gap refinements than a single residual bound. In contrast, the outer $c$ here is only globally Lipschitz within the Hadamard hypothesis above, so we rely on model decrease~\eqref{eq:pl-model-decrease} and inner steps satisfying templates~(DC)--(M) rather than on that convex-composite template; the exploration step of Section~\ref{sec:meta} still upgrades accumulation points to almost-sure d-stationarity under A1--A3, and greedy selection preserves the $O(1/\sqrt{N})$ scale on $\norm{x^k - z^{k+1}}$ from Lemma~\ref{lem:pl-descent} and Theorem~\ref{thm:pl-rate}.

% \begin{theorem}[Ergodic rate on the oracle step]\label{thm:pl-rate}
% Suppose $h^* > -\infty$ and, for each $k$, the oracle output satisfies~\eqref{eq:pl-model-decrease} so that Lemma~\ref{lem:pl-descent} applies.
% Then after $N$ steps of Algorithm~\ref{alg:random},
% \[
%   \min_{0 \leq k \leq N-1} \norm{x^k - z^{k+1}}
%   \;\leq\;
%   \sqrt{\frac{2\,\bigl(h(x^0) - h^*\bigr)}{(\rho - L_c L_F)\, N}}.
% \]
% \end{theorem}

% \begin{proof}
% See Subsection~\ref{subsec:app-pl-descent-and-rate} of Appendix~\ref{sec:appendix-pl}.
% \end{proof}

\subsection{Gradient Sampling Oracle and Goldstein Stationarity}\label{subsec:gs-oracle}

Consider the class $\mathcal{F}(\Delta, L)$ of $L$-Lipschitz, directionally differentiable, and bounded-below functions studied by \citet{Zhang2020}. 
The oracle is the Interpolated Normalized Gradient Descent (INGD) method of \citet{Zhang2020}, which at each iteration performs a local randomized search based on gradient interpolation and produces a feasible output $z^{k+1}$. 
This oracle requires only directional derivatives and function values, and does not rely on any DC decomposition or smooth structure.

% \paragraph{Compatibility with Assumptions A1--A3.}
Every $h \in \mathcal{F}(\Delta, L)$ is globally Lipschitz (A1). 
By construction, functions in this class are Hadamard directionally differentiable, which implies A2.
% , and Hadamard differentiability also yields continuity of $d \mapsto h'(x;\, d)$ for each fixed $x$ (A3).
Thus, Theorem~\ref{thm:random} applies directly to this setting.

\citet{Zhang2020} show that INGD finds a $(\delta, \varepsilon)$-Goldstein stationary point in 
$\tilde{O}(\Delta L^2 \varepsilon^{-3} \delta^{-1})$
oracle calls, where $\partial_\delta h(x) := \conv\{\nabla h(y) : y \in B(x,\delta)\}$ is the Goldstein $\delta$-subdifferential \citep{Goldstein1977}. 
Their analysis relies only on two properties: (i)~boundedness below of $h$, and (ii)~monotonic decrease along accepted steps.
Both are preserved under Algorithm~\ref{alg:random}: boundedness is unchanged, and greedy selection gives $h(x^{k+1}) \leq h(z^{k+1}) \leq h(x^k)$ independently of whether the exploration step is accepted.
Therefore, the $\tilde{O}(\Delta L^2 \varepsilon^{-3} \delta^{-1})$ complexity bound carries over without modification.

INGD targets Goldstein stationarity, a first-order condition based on approximate subdifferentials. 
The exploration step preserves this complexity guarantee while upgrading all accumulation points to almost-sure d-stationarity (Theorem~\ref{thm:random}). 
This demonstrates compatibility beyond structured settings, including purely derivative-based sampling schemes.

% INGD targets Goldstein stationarity, a first-order condition based on approximate subdifferentials. 
% The exploration step leaves this complexity guarantee intact while upgrading all accumulation points to almost-sure d-stationarity via Theorem~\ref{thm:random}. 
% This illustrates that the augmentation applies beyond structured settings such as DC or composite optimization, and is compatible with purely derivative-based sampling schemes.


% --------- Old Version --------
% \subsection{Gradient Sampling Oracle and Goldstein Stationarity}\label{subsec:gs-oracle}

% Consider the class $\mathcal{F}(\Delta, L)$ of $L$-Lipschitz, directionally differentiable, and bounded-below functions studied by \citet{Zhang2020}. The oracle is the Interpolated Normalized Gradient Descent (INGD) algorithm of \citet{Zhang2020}, which at each outer iteration $k$ performs a local search from $x^k$ using randomized gradient interpolation to produce a feasible output $z^{k+1}$. INGD requires access to directional derivatives and to function values; it does not rely on any DC decomposition or smooth structure.

% \paragraph{Verification of A1--A3.}
% Every $h \in \mathcal{F}(\Delta, L)$ is globally (hence locally) Lipschitz, so A1 holds. A2 holds by definition: $\mathcal{F}(\Delta, L)$ requires $h$ to be Hadamard directionally differentiable. A3 (continuity of $d \mapsto h'(x;\, d)$ for fixed $x$) follows from the Hadamard condition: Hadamard differentiability requires the joint limit $h'(x;\, d) = \lim_{(t,\,d') \to (0^+,\, d)} (h(x + td') - h(x))/t$ to hold along any path approaching $(0, d)$, not merely along the fixed ray $d' \equiv d$. Given $d^n \to d$, apply the joint limit at $(t_n, d^n)$ with $t_n \downarrow 0$ chosen so that $(h(x + t_n d^n) - h(x))/t_n$ approximates $h'(x;\, d^n)$; the joint-limit bound at direction $d$ then closes the triangle inequality, giving $h'(x;\, d^n) \to h'(x;\, d)$. In contrast, G\^{a}teaux differentiability (a limit in $t$ only, with $d$ fixed) carries no such uniformity in $d$ and does not imply A3.

% \paragraph{Rate preservation.}
% \citet{Zhang2020} prove that INGD finds a $(\delta, \varepsilon)$-Goldstein stationary point (a point $x$ with $\mathrm{dist}(0, \partial_\delta h(x)) \leq \varepsilon$) in $\tilde{O}(\Delta L^2 \varepsilon^{-3} \delta^{-1})$ oracle calls, where $\partial_\delta h(x) := \conv\{\nabla h(y) : y \in B(x,\delta)\}$ is the Goldstein $\delta$-subdifferential \citep{Goldstein1977}. Their proof uses only two properties of the iterates: (i)~$h$ is bounded below by $h^*$, and (ii)~$h(x^{k+1}) \leq h(x^k)$ whenever the descent condition fires. Both properties are preserved by Algorithm~\ref{alg:random}: property~(i) is unchanged, and property~(ii) continues to hold because greedy selection ensures $h(x^{k+1}) \leq h(z^{k+1}) \leq h(x^k)$. Therefore the $\tilde{O}(\Delta L^2 \varepsilon^{-3} \delta^{-1})$ complexity bound of \citet{Zhang2020} carries over to the augmented algorithm without modification.

% \input{experiments_summary}
% experiments_summary_story.tex
% Alternative body-section draft (story-first) for NeurIPS style.
% This file is an \input-able fragment without a preamble.

\section{Numerical Experiments}\label{sec:experiments-story}

Our experiments evaluate whether a single exploration step improves the terminal solution relative to the base oracle across representative benchmarks, using long runs to approach the asymptotic regime implied by the theory.
The full details are deferred to Appendix \ref{sec:appendix-experiments-v2}.

% Our experiments assess when one exploration step materially improves the terminal solution versus leaving it unchanged, using long runs to approach the asymptotic regime implied by the theory. We use the same augmentation on trimmed lasso (plain DCA often fails d-stationarity), least trimmed squares (DCA is reliable), and ReLU regression with a prox-linear oracle (non-DC).

\paragraph{Protocol in brief.}
For each instance, we compare the base oracle with its augmented version over three seeds, 
reporting $\Delta = h_{\text{base}} - \mathrm{median}_s\, h_{\text{aug},s}$ 
(positive values favor augmentation). \Cref{tab:experiments-story-summary} summarizes win/tie/loss counts, 
gap statistics, and d-stationarity failures.

% For each random instance we run the base oracle once and the augmented method on three seeds, scoring
% $\Delta = h_{\mathrm{base}} - \mathrm{median}_s\, h_{\mathrm{aug},s}$ (tie tolerance $10^{-12}$; positive $\Delta$ favors augmentation).
% \Cref{tab:experiments-story-summary} aggregates win/tie/loss counts, $\Delta$ statistics, McNemar tests, and d-stationarity diagnostic failures; setup and extra figures are in \Cref{sec:appendix-experiments-v2}, with diagnostic tolerances in \Cref{subsec:exp-dstat-checks}.
% Benchmark-specific d-stationarity checks are an exact DC subdifferential inclusion (trimmed lasso), an exact gradient match (LTS), and an exact DC subdifferential inclusion check for ReLU.



\paragraph{Trimmed lasso.}
In trimmed lasso, the difference between criticality and d-stationarity is evident in our diagnostics. In the focused setting
($m=50$, $n=100$, $k=5$, $\lambda=1$), plain DCA fails the d-stationarity
diagnostic on $99/100$ instances versus $4/100$ with augmentation, with objective outcomes $99/1/0$ wins/ties/losses under $\Delta$.
Directions are drawn from the gauss-axis mixture specified in \Cref{sec:appendix-experiments-v2}, which also records matched sphere runs under the same budgets -- uniform sphere sampling accepts virtually no exploration steps and behaves similarly to DCA.
The high base failure rate alongside the strong augmentation wins is consistent with a standard DC picture \citep{deOliveira2020}: where the nonsmooth term has nonsingleton subgradient, criticality ($\partial f(x)\cap\partial g(x)\neq\emptyset$) can hold without d-stationarity ($\partial g(x)\subseteq\partial f(x)$).

\paragraph{Least trimmed squares.}
Focused LTS ($m=100$, $n=50$, $q=10$): zero d-stationarity failures for base and augmented runs, $9/91/0$ wins/ties/losses, and a large win-conditional median $\Delta$ ($+52.2$). Augmentation is mostly inert but never hurts and occasionally finds a substantially better basin.

\paragraph{ReLU regression.}
Two focused ReLU settings ($200$ comparisons) with a prox-linear inner oracle and an exact DC subdifferential-inclusion d-stationarity check: zero d-stationarity failures, $44/156/0$ wins/ties/losses, mean $\Delta{=}+0.96$, McNemar $p{=}5.7\cdot 10^{-14}$, showing the same augmentation carries beyond DCA.

\begin{table}[t]
\centering
\caption{Headline outcomes across the three focused benchmarks.
\emph{w/t/l} denotes win/tie/loss under
$\Delta = h_{\mathrm{base}} - \mathrm{median}_s\, h_{\mathrm{aug},s}$.
\emph{\# non-d-stat (base / aug)} counts runs that fail the benchmark
d-stationarity diagnostic. ReLU aggregates two settings ($200$ paired
comparisons total).}
\label{tab:experiments-story-summary}
\small
\setlength{\tabcolsep}{4pt}
\begin{tabular}{l c c c c c c}
\toprule
Experiment & $n_{\mathrm{inst}}$ & w/t/l & mean $\Delta$ & med.\ win $\Delta$ & McNemar $p$ & \#~non-d-stat (base / aug) \\
\midrule
Trimmed lasso, $\lambda{=}1$        & 100 & 99/1/0    & $+0.29$ & $+0.27$ & $1.6 \cdot 10^{-30}$ & $99 / 4$ \\
LTS, $\sigma_{\mathrm{clean}}{=}4$  & 100 & 9/91/0    & $+3.96$ & $+52.2$ & $2.0 \cdot 10^{-3}$  & $0 / 0$  \\
ReLU regression                     & 200 & 44/156/0  & $+0.96$ & $+1.26$ & $5.7 \cdot 10^{-14}$ & $0 / 0$  \\
\bottomrule
\end{tabular}
\end{table}

% \begin{figure}[t]
% \centering
% \includegraphics[width=0.58\linewidth]{experiments/figures/fig_trajectory_gauss_axis.pdf}
% \caption{Representative trimmed-lasso trajectory in the failure-mode
% regime. Black: base DCA. Colored band/curve: min/median/max over
% augmentation seeds. The augmented method departs from the plateau reached
% by plain DCA.}
% \label{fig:experiments-story-trajectory}
% \end{figure}



\section{Conclusions and Limitations}\label{sec:limitations}\label{sec:conclusion}
\paragraph{Conclusions.} We showed that a single random feasible-direction step suffices to enforce almost-sure convergence of accumulation points to d-stationarity, independently of the underlying oracle. 
The key mechanism is a structural separation: the oracle determines convergence rates and practical performance, while exploration enforces asymptotic optimality at negligible cost. 
This separation is exact—rate guarantees of the base method are preserved via the inequality $h(x^{k+1}) \leq h(z^{k+1})$, while d-stationarity follows from the exploration step alone. 
The framework applies uniformly across oracle classes, including DCA, prox-linear methods, and gradient sampling, without modifying their analysis. 
These results suggest that lightweight directional exploration can serve as a modular tool for strengthening optimality guarantees. 
In particular, integrating structured feasible-direction mechanisms with model-based updates is a promising direction for advancing nonsmooth nonconvex optimization.

\paragraph{Limitations.} The guarantees are asymptotic. Theorem~\ref{thm:random} ensures almost-sure d-stationarity of accumulation points, but provides no finite-iteration bound for approaching or certifying it. The rates in Section~\ref{sec:oracles} concern the base oracle residual, not d-stationarity; thus finite-time proximity to a d-stationary point is not guaranteed.

Finite-time performance depends on the sampling distribution $\mathcal{D}_v$. Although the theory requires only full support on $S^{n-1} \cap L$, this may be ineffective in practice. In trimmed lasso, the uniform sphere sampler accepts virtually no exploration steps, whereas an axis-biased Gaussian mixture yields consistent gains by exploiting $\ell_1$-type sparsity. 
When such structure is unavailable and the base oracle is treated as a black box, designing an effective sampler is nontrivial, and uniform sphere sampling may be ineffective in practice.

% The guarantees established here are asymptotic. Theorem~\ref{thm:random} ensures almost-sure d-stationarity of accumulation points, but provides no finite-iteration bound for approaching or certifying d-stationarity. The rates in Section~\ref{sec:oracles} quantify convergence of the base oracle residual, not d-stationarity itself. Consequently, finite-time proximity to a d-stationary point is not guaranteed, even if the base oracle has converged under its own criterion.

% Finite-time performance is sensitive to the sampling distribution $\mathcal{D}_v$. The theory requires only full support on $S^{n-1} \cap L$, which need not yield effective exploration within a finite budget. This is evident in the trimmed lasso experiments: the uniform sphere sampler, while theoretically valid, yields virtually no accepted exploration steps in the focused regime and behaves comparably to plain DCA, whereas an axis-biased Gaussian mixture yields consistent improvements. The latter exploits the underlying $\ell_1$-type sparsity structure. When such structure is unavailable and the base oracle is treated as a black box, designing an effective sampler is nontrivial, and uniform sphere sampling may be ineffective in practice.

% \section{Limitations}\label{sec:limitations}

% The guarantees established in this work are asymptotic in nature. Theorem~\ref{thm:random} ensures that accumulation points are d-stationary almost surely, but provides no finite-iteration bound on when d-stationarity is approached or certified. The rate results of Section~\ref{sec:oracles} concern the convergence of the base oracle residual, not d-stationarity itself. In finite budgets, therefore, there is no guarantee that the iterates are close to a d-stationary point, even when the base oracle has converged by its own criterion.

% The finite-time behavior of the augmented method is strongly influenced by the choice of sampling distribution $\mathcal{D}_v$. While the theory requires only full support on $S^{n-1} \cap L$, this condition is not always sufficient for practical effectiveness in finite budgets. The trimmed lasso experiments illustrate this concretely: the uniform sphere sampler, despite satisfying all theoretical requirements, accepts essentially no exploration moves in the focused regime and is indistinguishable from plain DCA, while the axis-biased Gaussian mixture sampler produces substantial improvements. The sampler design used there was informed by the known $\ell_1$-type coordinate sparsity structure. When the base oracle is treated as a black box and problem structure is unavailable, designing an effective sampler becomes nontrivial, and uniform sphere sampling may be a poor default.

% \section{Conclusion}\label{sec:conclusion}

% We showed that a single random feasible-direction step suffices to enforce almost-sure convergence of accumulation points to d-stationarity, independently of the underlying oracle. 
% The key mechanism is a structural separation: the oracle determines convergence rates and practical performance, while exploration enforces asymptotic optimality at negligible cost. 
% This separation is exact—rate guarantees of the base method are preserved via the inequality $h(x^{k+1}) \leq h(z^{k+1})$, while d-stationarity follows from the exploration step alone. 
% The framework applies uniformly across oracle classes, including DCA, prox-linear methods, and gradient sampling, without modifying their analysis. 
% These results suggest that lightweight directional exploration can serve as a modular tool for strengthening optimality guarantees. 
% In particular, integrating structured feasible-direction mechanisms with model-based updates is a promising direction for advancing nonsmooth nonconvex optimization.

% \section{Conclusion}\label{sec:conclusion}

% We have introduced a meta-algorithm that guarantees almost-sure asymptotic convergence to d-stationary points of a nonsmooth nonconvex objective over a general closed convex set, regardless of the inner oracle. The key insight is a structural separation: the oracle drives rates and problem-specific performance, while the exploration step, which samples both a feasible direction and a step size uniformly and accepts the result via a two-point comparison, carries the entire asymptotic guarantee at a cost of one feasibility check and at most one function evaluation per iteration.

% Concrete instantiations with the DCA oracle (Section~\ref{sec:dc}) and with prox-linear-type oracles (Section~\ref{subsec:pl-oracle}, Theorem~\ref{thm:pl-rate}) demonstrate that the meta-algorithm framework is rate-compatible: in each case an $O(1/\sqrt{N})$ scale bound on the oracle residual $\min_k \norm{x^k - z^{k+1}}$ is preserved exactly, via the single inequality $h(x^{k+1}) \leq h(z^{k+1})$ that greedy selection provides. The exploration step can only improve the objective constant. Accumulation points inherit the base-oracle stationarity interpretation of that residual (where applicable) and are additionally d-stationary almost surely (Theorem~\ref{thm:random}), with these quantitative and geometric guarantees arising from structurally independent parts of the analysis.

% \paragraph{Open questions and future directions.}
% \begin{itemize}
% \item \textit{Rates beyond the DCA, prox-linear, and INGD cases.} Can a convergence rate be established for the meta-algorithm beyond the $O(1/\sqrt{N})$ rates on the oracle residual inherited from the DCA and prox-linear oracles of Sections~\ref{sec:dc}--\ref{sec:oracles}?
% \item \textit{Stochastic objectives.} Can the framework be extended to stochastic objectives where $h$ is observed only via noisy function evaluations?
% \item \textit{Beyond DC structure.} Assumptions~A1--A3 require only existence and continuity of directional derivatives; the DC structure is not used in the proofs. Which broader classes of objectives can be handled?
% \item \textit{Practical oracle design.} Identifying specific oracles that, when augmented with the exploration step, yield both fast rates and d-stationarity for structured machine learning problems is an important direction.
% \end{itemize}

\begin{ack}
No external funding was received for this work. The authors declare no competing financial interests.
\end{ack}

\bibliographystyle{plainnat}
\bibliography{references}

\appendix
% appendix_convergence_proofs.tex
% \input-able fragment (no preamble). Deferred proofs for Section~\ref{sec:convergence}.
\section{Proofs of \Cref{sec:convergence}}
\label{sec:appendix Proofs of}

% \begin{proof}[Proof of Lemma \ref{lem:BC-directions}]
% We assume throughout that $\dim L \ge 1$; otherwise $C$ is a singleton and the result is trivial.

% Let $\mathcal{F}_{k-1}$ denote the $\sigma$-algebra generated by $\{x^0, v^0_{ex}, \hat{t}^0, \ldots, x^{k-1}, v^{k-1}_{ex}, \hat{t}^{k-1}, x^k\}$, so that $x^k$ is $\mathcal{F}_{k-1}$-measurable while both $v^k_{ex}$ and $\hat{t}^k$ are independent of $\mathcal{F}_{k-1}$ (and of each other).

% Fix a point $y \in C$, a radius $\delta > 0$, a spherical cap
% \[
% V := B[w, \varepsilon'] \cap S^{n-1} \cap L,
% \]
% for some $w \in S^{n-1} \cap L$ and $\varepsilon' > 0$, and an interval $[a', b'] \subseteq (0, r]$ with $a' < b'$. 
% Set
% \[
% p := \mathbb{P}(v^k_{ex} \in V) > 0,\qquad
% q := \mathbb{P}(\hat t^k \in [a', b']) > 0,
% \]
% where positivity follows from the full-support assumptions.

% Define
% \[
% A_k := \{v^k_{ex} \in V\} \cap \{\hat{t}^k \in [a', b']\} \cap \{\norm{x^k - y} < \delta\}.
% \]
% Independence yields
% \[
% \mathbb{P}(A_k \mid \mathcal{F}_{k-1})
% = \mathbf{1}_{\{\norm{x^k - y} < \delta\}} \cdot p \cdot q.
% \]
% By the conditional Borel--Cantelli--L\'evy lemma \citep[Ch.~VII, \S5, Cor.~2]{Shiryaev2019},
% \begin{equation}\label{eq:BCL-step}
%   \Bigl\{\sum_{k} \mathbb{P}(A_k \mid \mathcal{F}_{k-1}) = \infty\Bigr\}
%   =
%   \{A_k\ \text{occurs infinitely often}\}
%   \quad\text{a.s.}
% \end{equation}
% In particular, whenever $x^k$ visits $B(y,\delta)$ infinitely often, $A_k$ occurs infinitely often almost surely.

% Let $\{y_n\}$ be a countable dense subset of $C$, $\{w_i\}$ a countable dense subset of $S^{n-1} \cap L$, and let $(\delta_m, \varepsilon'_l, a'_s, b'_s)$ range over positive rationals with $a'_s < b'_s \leq r$. Applying~\eqref{eq:BCL-step} to each tuple and taking a countable intersection of probability-$1$ events, we obtain a full-probability event $\Omega_0$ on which the following grid property holds simultaneously:

% For every tuple $(y_n, \delta_m, w_i, \varepsilon'_l, a'_s, b'_s)$, if $\norm{x^k - y_n} < \delta_m$ for infinitely many $k$, then there are infinitely many such $k$ with additionally
% \[
% v^k_{ex} \in B[w_i, \varepsilon'_l] \cap S^{n-1} \cap L,
% \qquad
% \hat{t}^k \in [a'_s, b'_s].
% \]

% Lets us consider $\Omega_0$. 
% Let $x^*$ be any accumulation point of $\{x^k\}$ and fix $v \in \ri(D_C(x^*)) \cap S^{n-1}$. Fix arbitrary $\eta \in (0,\rho)$. By density, choose rationals $y_n, \delta_m, w_i, \varepsilon'_l, a'_s, b'_s$ such that
% $
% \|y_n - x^*\| < \eta/4,
% \delta_m \in (\eta/2, 3\eta/4),
% \|w_i - v\| < \eta/4,
% \varepsilon'_l < \eta/4,
% $
% and
% $
% a < a'_s < b'_s < b.
% $
% Since $x^*$ is an accumulation point, $\|x^k - x^*\| < \eta/4$ for infinitely many $k$. For such indices,
% \[
% \|x^k - y_n\|
% \le \|x^k - x^*\| + \|x^* - y_n\|
% < \eta/2 < \delta_m.
% \]
% Hence, by the grid property, there are infinitely many indices $k$ such that
% \[
% \|x^k - y_n\| < \delta_m,\quad
% v^k_{ex} \in B[w_i, \varepsilon'_l],\quad
% \hat{t}^k \in [a'_s, b'_s].
% \]

% For these indices,
% \[
% \|x^k - x^*\|
% \le \|x^k - y_n\| + \|y_n - x^*\|
% < \delta_m + \eta/4
% < \eta,
% \]
% and
% \[
% \|v^k_{ex} - v\|
% \le \|v^k_{ex} - w_i\| + \|w_i - v\|
% < \varepsilon'_l + \eta/4
% < \eta < \rho,
% \]
% while $\hat{t}^k \in [a'_s, b'_s] \subseteq (a,b)$.

% Let $\eta_q \downarrow 0$ with $\eta_1 < \rho$, and for each $q$ let $K_q$ be the infinite set of indices satisfying the above properties with $\eta=\eta_q$. Choose inductively $j_q \in K_q$ with $j_q > j_{q-1}$. Then, since $\|x^{j_q} - x^*\| < \eta_q$ and $\eta_q \to 0$, we have $x^{j_q} \to x^*$. Similarly, $\|v^{j_q}_{ex} - v\| < \eta_q$ implies $v^{j_q}_{ex} \to v$, while $\|v^{j_q}_{ex} - v\| < \rho$ and $\hat{t}^{j_q} \in [a,b]$ for all $q$. This completes the proof.
% \end{proof}

% \medskip

\begin{proof}[Proof of Lemma \ref{lem:step-bound}]
Let $k := \dim(L)$. Since $v \in \ri(D_C(x^*))$, there exists $r_0 > 0$ such that
\[
B(v, r_0) \cap L \subseteq \ri(D_C(x^*)).
\]
Pick $k+1$ directions $d_0, \ldots, d_k \in B(v, r_0) \cap L$ whose convex hull is a $k$-dimensional simplex in $L$ containing $v$ in its relative interior. Then
\[
\conv\{d_0, \ldots, d_k\} \subseteq \ri(D_C(x^*)).
\]
Since each $d_j \in D_C(x^*)$, choose $\varepsilon_j > 0$ with $x^* + \varepsilon_j d_j \in C$, and set $\bar{\varepsilon} := \min_j \varepsilon_j > 0$. By convexity of $C$, $p_j := x^* + \bar{\varepsilon}\, d_j \in C$ for each $j$, and for any convex combination $d = \sum_j \lambda_j d_j$,
\[
  x^* + \bar{\varepsilon}\, d = \sum_{j} \lambda_j\, p_j \in C.
\]
Since $v$ lies in the relative interior of the simplex, there exists $\rho_0 > 0$ with
\[
B(v, \rho_0) \cap L \subseteq \conv\{d_0, \ldots, d_k\},
\]
and hence $x^* + \bar{\varepsilon}\, w \in C$ for every $w \in B(v, \rho_0) \cap L$. The map $w \mapsto x^* + \bar{\varepsilon}\, w$ sends $B(v,\rho_0)\cap L$ onto
\[
B(x^* + \bar{\varepsilon}\, v,\; \bar{\varepsilon}\rho_0)\cap \aff(C),
\]
and therefore
\[
  B\!\left(x^* + \bar{\varepsilon}\, v,\; \bar{\varepsilon}\rho_0\right) \cap \aff(C) \subseteq C.
\]

Set $\rho := \rho_0 / 2$ and $\delta := \bar{\varepsilon}\rho_0 / 2$. Then (i) holds because
\[
B(v,\, \rho) \cap L \subseteq B(v,\, \rho_0) \cap L \subseteq \conv\{d_0, \ldots, d_k\} \subseteq \ri(D_C(x^*)).
\]
For (ii), fix $x \in C$ with $\norm{x - x^*} < \delta$ and $w \in B(v,\, \rho) \cap L$. Then
\[
  \norm{(x + \bar{\varepsilon}\, w) - (x^* + \bar{\varepsilon}\, v)} \leq \norm{x - x^*} + \bar{\varepsilon}\norm{w - v} < \delta + \bar{\varepsilon}\rho = \bar{\varepsilon}\rho_0,
\]
and since $x \in C \subseteq \aff(C)$ and $w \in L$, we have $x + \bar{\varepsilon}\, w \in \aff(C)$; hence
\[
x + \bar{\varepsilon}\, w \in B(x^* + \bar{\varepsilon}\, v,\, \bar{\varepsilon}\rho_0) \cap \aff(C) \subseteq C.
\]
Convexity then gives
\[
  x + t w = (1 - t/\bar{\varepsilon})\, x + (t/\bar{\varepsilon})\,(x + \bar{\varepsilon}\, w) \in C \qquad \text{for every } t \in [0,\, \bar{\varepsilon}],
\]
so $\tau(x;\, w) \geq \bar{\varepsilon}$.
\end{proof}

\medskip

\begin{proof}[Proof of Lemma \ref{lem:continuity}]
Let $K := D_C(x^*)$, a convex cone. 
By Lemma~\ref{ass:A3}, the map $d \mapsto h'(x^*;\, d)$ is continuous.
We next show positive homogeneity. For any $\alpha > 0$,
\[
  h'(x^*;\, \alpha d)
  = \lim_{t\downarrow 0}\frac{h(x^* + t\alpha d)-h(x^*)}{t}
  = \alpha \lim_{t\downarrow 0}\frac{h(x^* + t d)-h(x^*)}{t}
  = \alpha h'(x^*;\, d).
\]

For the extension statement, recall that for any convex set $K$,
\[
\cl(\ri(K)) = \cl(K),
\]
where the relative interior is taken with respect to $\aff(K)=L$. Let $d \in \cl(K)$. Then there exists a sequence $v_m \in \ri(K)$ with $v_m \to d$. By continuity,
\[
  h'(x^*;\, d) = \lim_{m\to\infty} h'(x^*;\, v_m) \geq 0,
\]
since $h'(x^*;\, v_m) \geq 0$ for all $m$.
\end{proof}


% appendix_dca.tex
% \input-able fragment (no preamble). Companion to Section~\ref{sec:dc}.
\section{Proofs for the DCA oracle rate (Section~\ref{sec:dc})}
\label{sec:appendix-dca}

\begin{lemma}[Descent at the oracle output]\label{lem:dca-descent}
For each $k$, $h(x^k) - h(z^{k+1}) \geq \frac{\mu_1 + \mu_2}{2}\norm{x^k - z^{k+1}}^2$.
\end{lemma}

\begin{proof}
By $\mu_1$-strong convexity of $f_1$ at the subgradient $s^{k+1} \in \partial f_1(z^{k+1})$:
\[
  f_1(x^k) \geq f_1(z^{k+1}) + \inner{s^{k+1}}{x^k - z^{k+1}} + \tfrac{\mu_1}{2}\norm{x^k - z^{k+1}}^2.
\]
First-order optimality of the DCA subproblem at $z^{k+1}$ gives $s^{k+1} \in \partial f_1(z^{k+1})$, $n^{k+1} \in N_C(z^{k+1})$, and
\begin{equation}\label{eq:OC}
  s^{k+1} + n^{k+1} = g^k_2.
\end{equation}
Substituting $s^{k+1} = g^k_2 - n^{k+1}$ from~\eqref{eq:OC} and using $\inner{n^{k+1}}{x^k - z^{k+1}} \leq 0$ (since $n^{k+1} \in N_C(z^{k+1})$ and $x^k \in C$):
\begin{equation}\label{eq:f1-simplified}
  f_1(x^k) \geq f_1(z^{k+1}) + \inner{g^k_2}{x^k - z^{k+1}} + \tfrac{\mu_1}{2}\norm{x^k - z^{k+1}}^2.
\end{equation}
By $\mu_2$-strong convexity of $f_2$ with subgradient $g^k_2 \in \partial f_2(x^k)$:
\begin{equation}\label{eq:f2-bound}
  f_2(x^k) \leq f_2(z^{k+1}) + \inner{g^k_2}{x^k - z^{k+1}} - \tfrac{\mu_2}{2}\norm{x^k - z^{k+1}}^2.
\end{equation}
Subtracting~\eqref{eq:f2-bound} from~\eqref{eq:f1-simplified}, the $\inner{g^k_2}{\cdot}$ terms cancel:
$h(x^k) - h(z^{k+1}) \geq \frac{\mu_1 + \mu_2}{2}\norm{x^k - z^{k+1}}^2$.
\end{proof}

\begin{proof}[Proof of Theorem~\ref{thm:dca-rate}]
By greedy selection, $h(x^{k+1}) \leq h(z^{k+1})$. Combining with Lemma~\ref{lem:dca-descent}:
\[
  h(x^k) - h(x^{k+1}) \geq \frac{\mu_1 + \mu_2}{2}\norm{x^k - z^{k+1}}^2.
\]
Summing over $k = 0, \ldots, N-1$ and telescoping:
\[
  \frac{\mu_1 + \mu_2}{2}\sum_{k=0}^{N-1}\norm{x^k - z^{k+1}}^2 \leq h(x^0) - h(x^N) \leq h(x^0) - h^*.
\]
Taking the minimum over $k$ (at most the average) and the square root completes the proof.
\end{proof}

\begin{remark}[Oracle residual $\norm{x^k - z^{k+1}}$]\label{rem:pl-oracle-residual}\label{rem:dca-optimality-measure}
For the DCA oracle, $\norm{x^k - z^{k+1}} = 0$ if and only if $x^k$ is critical (a fixed point of the oracle subproblem). Thus $\norm{x^k - z^{k+1}}$ is a quantitative residual for distance to criticality, and Theorem~\ref{thm:dca-rate} gives the preserved ergodic $O(1/\sqrt{N})$ rate in exactly this residual.

This rate statement concerns a finite-iteration optimality measure, whereas the asymptotic guarantee of the augmented method (Theorem~\ref{thm:random}) is almost-sure d-stationarity of accumulation points. There is no conflict: the base algorithm converges to criticality, while augmentation upgrades the limit notion to d-stationarity, which is strictly stronger in general (Remark~\ref{rem:stationarity-comparison}).
\end{remark}
% appendix_prox_linear_stationarity.tex
% \input-able fragment (no preamble). Companion to Section~\ref{subsec:pl-oracle}.

\section{Prox-linear-type oracles: model decrease, descent, and limiting stationarity}
\label{sec:appendix-pl}

This appendix supports Section~\ref{subsec:pl-oracle}.
Throughout, $z^{k+1}$ denotes the candidate produced by a prox-linear-type oracle at the current iterate $x^k$ of Algorithm~\ref{alg:random} (unconstrained $C=\R^n$ in~\eqref{eq:prox-linear}).
We write $\ell^k(y) := F(x^k) + \nabla F(x^k)(y-x^k)$ and $m^k(y) := c(\ell^k(y)) + \tfrac{\rho}{2}\norm{y-x^k}^2$.
The descent calculation is classical for Lipschitz composite models \citep{Drusvyatskiy2019}.

\subsection{Lemma~\ref{lem:pl-descent} and Theorem~\ref{thm:pl-rate}}\label{subsec:app-pl-descent-and-rate}

\begin{lemma}[Quadratic descent from model decrease]\label{lem:pl-descent}
Under the Lipschitz assumptions on $c$ and $F$ and with $\rho > L_c L_F$, if~\eqref{eq:pl-model-decrease} holds at iteration $k$, then
\begin{equation}\label{eq:pl-qdescent-state}
  h(x^k) - h(z^{k+1}) \;\geq\; \tfrac{1}{2}(\rho - L_c L_F)\,\norm{z^{k+1} - x^k}^2.
\end{equation}
\end{lemma}

\begin{proof}
From $m^k(z^{k+1}) \leq m^k(x^k) = h(x^k)$,
\begin{equation}\label{eq:app-pl-step1}
  h(x^k) - c(\ell^k(z^{k+1})) \;\geq\; \tfrac{\rho}{2}\norm{z^{k+1}-x^k}^2.
\end{equation}
By the Lipschitz Jacobian estimate for $F$ \citep{Drusvyatskiy2019},
$\norm{F(z^{k+1}) - \ell^k(z^{k+1})} \leq \tfrac{L_F}{2}\norm{z^{k+1}-x^k}^2$, and Lipschitzness of $c$ gives
$c(\ell^k(z^{k+1})) \geq c(F(z^{k+1})) - L_c \norm{F(z^{k+1})-\ell^k(z^{k+1})}
\geq h(z^{k+1}) - \tfrac{L_c L_F}{2}\norm{z^{k+1}-x^k}^2$.
Substituting into~\eqref{eq:app-pl-step1} yields~\eqref{eq:pl-qdescent-state}.
If moreover $h$ is bounded below and $h(x^{k+1}) \leq h(z^{k+1})$ at every $k$, then summing the inequality of Lemma~\ref{lem:pl-descent} gives $\sum_k \norm{z^{k+1}-x^k}^2 < \infty$ and hence $\norm{z^{k+1}-x^k}\to 0$.
\end{proof}

\begin{proof}[Proof of Theorem~\ref{thm:pl-rate}]
Greedy selection gives $h(x^{k+1}) \leq h(z^{k+1})$. Combining with~\eqref{eq:pl-qdescent-state} yields
\begin{equation}
  h(x^k) - h(z^{k+1}) \geq \alpha\norm{x^k - z^{k+1}}^2 \tag{QD}\label{eq:QD}
\end{equation}
 with $\alpha = \tfrac{1}{2}(\rho - L_c L_F)$. Telescoping the same way as in Theorem~\ref{thm:dca-rate} proves the bound.
\end{proof}

\begin{assumption}[PL-BDD]\label{ass:app-pl-bdd}
The sublevel set $\{x \in \R^n : h(x) \leq h(x^0)\}$ is bounded.
\end{assumption}

\begin{remark}
Under greedy selection, $h(x^{k+1}) \leq h(x^k)$ for all $k$, so Assumption~\ref{ass:app-pl-bdd} implies that $\{x^k\}$ is bounded and admits accumulation points.
This assumption is used only to guarantee nonempty accumulation-point limits in the results below; the limit identities along any convergent subsequence do not otherwise require global boundedness of $\{x^k\}$.
\end{remark}



\paragraph{Reading guide for the next theorems.}\label{rem:app-pl-purpose}
The next three theorems identify what vanishing prox-linear residual
$\norm{z^{k+1}-x^k}$ implies at accumulation points under three inner-termination templates: DC, Clarke, and Mordukhovich.
This is the same residual that appears in Lemma~\ref{lem:pl-descent} and Theorem~\ref{thm:pl-rate}, so the ergodic $O(1/\sqrt{N})$ bound is a rate statement in a concrete optimality metric.
In the augmented algorithm, accumulation points are upgraded further to almost-sure d-stationarity (Theorem~\ref{thm:random}), which is strictly stronger in general (Remark~\ref{rem:stationarity-comparison}; see also Remark~\ref{rem:pl-oracle-residual}).

\subsection{DC outer function}

\begin{assumption}[PL-DC1]\label{ass:app-pl-dc-c}
The outer function $c \colon \R^m \to \R$ is DC: $c(y)=f(y)-g(y)$ with $f,g \colon \R^m \to \R$ convex and finite-valued, and $c$ is globally Lipschitz with constant $L_c>0$.
\end{assumption}

\begin{assumption}[PL-DC2]\label{ass:app-pl-dc-F}
The inner map $F \colon \R^n \to \R^m$ is continuously differentiable and $\nabla F$ is $L_F$-Lipschitz in operator norm, as in Lemma~\ref{lem:pl-descent}.
\end{assumption}

\begin{assumption}[PL-DC3]\label{ass:app-pl-dc-rho}
The parameter satisfies $\rho > L_c L_F$.
\end{assumption}

\paragraph{Inner termination (DC criticality of the model).}
Fix $k$ and decompose $m^k = \varphi^k - \psi^k$ with
$\varphi^k(y) := f(\ell^k(y)) + \tfrac{\rho}{2}\norm{y-x^k}^2$ and $\psi^k(y) := g(\ell^k(y))$ (both convex).
The oracle output $z^{k+1}$ is \emph{DC critical} for $m^k$ if
\begin{equation}\label{eq:app-pl-dc-inner}
  \nabla F(x^k)^\top \xi^{k+1} + \rho(z^{k+1}-x^k) = \nabla F(x^k)^\top \eta^{k+1}
\end{equation}
for some $\xi^{k+1} \in \partial f(\ell^k(z^{k+1}))$ and $\eta^{k+1} \in \partial g(\ell^k(z^{k+1}))$, and satisfies~\eqref{eq:pl-model-decrease}.

\begin{theorem}[DC criticality at accumulation points]\label{thm:app-pl-dc}
Under Assumptions~\ref{ass:app-pl-dc-c}--\ref{ass:app-pl-dc-rho} and~\ref{ass:app-pl-bdd}, suppose each oracle output $z^{k+1}$ satisfies~\eqref{eq:app-pl-dc-inner} and~\eqref{eq:pl-model-decrease}.
Then every accumulation point $x^*$ of $\{x^k\}$ satisfies
\begin{equation}\label{eq:app-pl-dc-stat}
  \nabla F(x^*)^\top \partial f(F(x^*)) \;\cap\; \nabla F(x^*)^\top \partial g(F(x^*)) \;\neq\; \emptyset,
\end{equation}
i.e., DC criticality of $h=f\circ F - g\circ F$ at $x^*$.
\end{theorem}

\begin{proof}
By Lemma~\ref{lem:pl-descent} and greedy selection, $\sum_k \norm{z^{k+1}-x^k}^2 < \infty$, hence $\norm{z^{k+1}-x^k}\to 0$.
Also $h(x^{k+1}) \leq h(x^k)$, so $h(x^k) \leq h(x^0)$ for all $k$ and $\{x^k\}$ lies in a compact set by Assumption~\ref{ass:app-pl-bdd}; thus $\{x^k\}$ has accumulation points.

Let $x^{k_j} \to x^*$. Since $\norm{z^{k_j+1}-x^{k_j}}\to 0$, also $z^{k_j+1} \to x^*$.
Along $k_j$, $\nabla F(x^{k_j}) \to \nabla F(x^*)$.
Set $y^j := \ell^{k_j}(z^{k_j+1}) = F(x^{k_j}) + \nabla F(x^{k_j})(z^{k_j+1}-x^{k_j}) \to F(x^*)$.
For $\xi^{k_j+1} \in \partial f(y^j)$ and $\eta^{k_j+1} \in \partial g(y^j)$, the sequences $\{\xi^{k_j+1}\}$ and $\{\eta^{k_j+1}\}$ are bounded: for all large $j$, $y^j$ lies in a fixed compact set, and convex finite $f,g$ are Lipschitz on compacts, hence their subdifferentials are uniformly bounded there.
Passing to a subsequence, $\xi^{k_j+1} \to \xi^*$ and $\eta^{k_j+1} \to \eta^*$, and outer semicontinuity of convex subdifferentials \citep[Theorem~24.4]{Rockafellar1997} yields $\xi^* \in \partial f(F(x^*))$ and $\eta^* \in \partial g(F(x^*))$.
The proximal term $\rho(z^{k_j+1}-x^{k_j}) \to 0$.
Passing to the limit in~\eqref{eq:app-pl-dc-inner} gives $\nabla F(x^*)^\top \xi^* = \nabla F(x^*)^\top \eta^*$, which is~\eqref{eq:app-pl-dc-stat}.
\end{proof}

% \begin{remark}
% For affine $\ell^k$, the subdifferential chain rule $\partial(f\circ \ell^k)(y) = \nabla F(x^k)^\top \partial f(\ell^k(y))$ holds without surjectivity \citep[Theorem~23.9]{Rockafellar1997}, which justifies~\eqref{eq:app-pl-dc-inner}.
% \end{remark}

\subsection{Clarke outer function}

\begin{assumption}[PL-C1]\label{ass:app-pl-clarke-c}
The function $c \colon \R^m \to \R$ is globally Lipschitz with constant $L_c>0$ (hence locally Lipschitz and therefore Clarke subdifferentiable).
\end{assumption}

\begin{assumption}[PL-C2]\label{ass:app-pl-clarke-F}
The map $F$ satisfies the same smoothness as in Lemma~\ref{lem:pl-descent}, and $\nabla F(x)$ is surjective (full row rank $m$) for every $x \in \R^n$.
\end{assumption}

\begin{assumption}[PL-C3]\label{ass:app-pl-clarke-rho}
$\rho > L_c L_F$, and $h=c\circ F$ is bounded below on $\R^n$ (equivalently, $h^* > -\infty$ in the notation of Section~\ref{subsec:pl-oracle}).
\end{assumption}

\paragraph{Inner termination (Clarke stationarity of the model).}
The oracle output $z^{k+1}$ satisfies
\begin{equation}\label{eq:app-pl-clarke-inner}
  0 \in \partialC\!\Bigl( c \circ \ell^k + \tfrac{\rho}{2}\norm{\,\cdot - x^k}^2 \Bigr)(z^{k+1}),
\end{equation}
together with~\eqref{eq:pl-model-decrease}.
Under surjectivity of $\nabla F(x^k)$, the Clarke chain rule \citep[Theorem~2.3.10]{clarke1990optimization} and sum rule \citep[Proposition~2.3.3]{clarke1990optimization} expand~\eqref{eq:app-pl-clarke-inner} to
\begin{equation}\label{eq:app-pl-clarke-expanded}
  0 \in \nabla F(x^k)^\top \partialC c(\ell^k(z^{k+1})) + \rho(z^{k+1}-x^k).
\end{equation}

\begin{theorem}[Clarke stationarity at accumulation points]\label{thm:app-pl-clarke}
Under Assumptions~\ref{ass:app-pl-clarke-c}--\ref{ass:app-pl-clarke-rho} and~\ref{ass:app-pl-bdd}, suppose each $z^{k+1}$ satisfies~\eqref{eq:app-pl-clarke-inner} and~\eqref{eq:pl-model-decrease}.
Then every accumulation point $x^*$ of $\{x^k\}$ satisfies $0 \in \partialC h(x^*)$, i.e., Clarke stationarity of $h=c\circ F$.
\end{theorem}

\begin{proof}
Greedy monotonicity and Assumption~\ref{ass:app-pl-bdd} imply that $\{x^k\}$ is bounded, hence it has accumulation points.
Let $x^{k_j} \to x^*$.
Lemma~\ref{lem:pl-descent} and greedy selection give $\norm{z^{k_j+1}-x^{k_j}}\to 0$, hence $z^{k_j+1} \to x^*$.
Choose $v^{k_j+1} \in \partialC c(\ell^{k_j}(z^{k_j+1}))$ with $\nabla F(x^{k_j})^\top v^{k_j+1} + \rho(z^{k_j+1}-x^{k_j})=0$.
The Clarke subdifferential is locally bounded \citep[Proposition~2.1.2(a)]{clarke1990optimization}; pass to a subsequence with $v^{k_j+1} \to v^*$, and outer semicontinuity of $\partialC c$ yields $v^* \in \partialC c(F(x^*))$.
Taking limits, $\nabla F(x^*)^\top v^* = 0$, i.e., $0 \in \nabla F(x^*)^\top \partialC c(F(x^*))$.
Under surjectivity, the Clarke chain rule for $c\circ F$ \citep[Theorem~2.3.10]{clarke1990optimization} gives $\partialC h(x^*) = \nabla F(x^*)^\top \partialC c(F(x^*))$ (equality), hence $0 \in \partialC h(x^*)$.
\end{proof}

\subsection{Mordukhovich (limiting) outer function}

\paragraph{Inner termination (Mordukhovich stationarity of the model).}
Replace~\eqref{eq:app-pl-clarke-inner} by
\begin{equation}\label{eq:app-pl-mord-inner}
  0 \in \partialL\!\Bigl( c \circ \ell^k + \tfrac{\rho}{2}\norm{\,\cdot - x^k}^2 \Bigr)(z^{k+1}),
\end{equation}
with~\eqref{eq:pl-model-decrease}, where $\partialL$ denotes the limiting (Mordukhovich) subdifferential.
Under surjectivity of $\nabla F(x^k)$, the limiting chain rule \citep[Exercise~10.7]{RockafellarWets2009} and the smooth-sum rule for $\partialL$ \citep[Exercise~10.10]{RockafellarWets2009} yield the expansion
\begin{equation}\label{eq:app-pl-mord-expanded}
  0 \in \nabla F(x^k)^\top \partialL c(\ell^k(z^{k+1})) + \rho(z^{k+1}-x^k).
\end{equation}

\begin{theorem}[Mordukhovich stationarity at accumulation points]\label{thm:app-pl-mord}
Under Assumptions~\ref{ass:app-pl-clarke-c}--\ref{ass:app-pl-clarke-rho} and~\ref{ass:app-pl-bdd}, suppose each $z^{k+1}$ satisfies~\eqref{eq:app-pl-mord-inner} and~\eqref{eq:pl-model-decrease}.
Then every accumulation point $x^*$ of $\{x^k\}$ satisfies $0 \in \partialL h(x^*)$, i.e., Mordukhovich stationarity of $h=c\circ F$.
\end{theorem}

\begin{proof}
Greedy monotonicity and Assumption~\ref{ass:app-pl-bdd} imply that $\{x^k\}$ is bounded, hence it has accumulation points.
Let $x^{k_j} \to x^*$; by Lemma~\ref{lem:pl-descent} and greedy selection, $\norm{z^{k_j+1}-x^{k_j}}\to 0$, hence $z^{k_j+1} \to x^*$.
Pick $v^{k_j+1} \in \partialL c(\ell^{k_j}(z^{k_j+1}))$ with $\nabla F(x^{k_j})^\top v^{k_j+1} + \rho(z^{k_j+1}-x^{k_j})=0$.
Since $\partialL c(y) \subseteq \partialC c(y)$ for locally Lipschitz $c$ \citep[Theorems~8.49 and~9.13]{RockafellarWets2009} and $\partialC c$ is locally bounded, the $v^{k_j+1}$ are bounded; pass to a convergent subsequence $v^{k_j+1} \to v^*$.
Outer semicontinuity of $\partialL c$ gives $v^* \in \partialL c(F(x^*))$, and taking limits yields $0 \in \nabla F(x^*)^\top \partialL c(F(x^*))$.
Under surjectivity, $\partialL h(x^*) = \nabla F(x^*)^\top \partialL c(F(x^*))$ \citep[Exercise~10.7]{RockafellarWets2009}, hence $0 \in \partialL h(x^*)$.
\end{proof}

\begin{remark}[Comparison]\label{rem:app-pl-compare}
The descent step is common (Lemma~\ref{lem:pl-descent}); only the subdifferential calculus in the limit and the final chain rule differ.
Since $\partialL c(y) \subseteq \partialC c(y)$ always, Mordukhovich stationarity of $h$ is stronger than Clarke stationarity, and~\eqref{eq:app-pl-mord-inner} is a stronger inner requirement than~\eqref{eq:app-pl-clarke-inner}.
Assumption~\ref{ass:app-pl-bdd} is imposed uniformly in Theorems~\ref{thm:app-pl-dc}--\ref{thm:app-pl-mord} so that greedy iterates admit accumulation points; it is not needed for the algebraic form of the limit along a convergent subsequence.
\end{remark}

% experiments_appendix_v2.tex
% Narrative-style appendix variant for experiments.
% This file is an \input-able fragment without a preamble.

\section{Additional Experimental Details and Discussion}
\label{sec:appendix-experiments-v2}

The purpose of this appendix is to document, in one place, when the exploration step
changes the outcome and whether those changes are favorable. Across all experiments, we
compare a deterministic base oracle to its augmented
version (oracle step plus one randomized exploration step), using the same initialization
per instance and three augmentation seeds. We report win/tie/loss against the base
solution using the signed gap
\[
  \Delta \;=\; h_{\mathrm{base}} - \mathrm{median}_s\, h_{\mathrm{aug},s},
\]
with tolerance $10^{-12}$. Positive $\Delta$ means augmentation is better. In addition to
mean and median gap (wins), we report McNemar's exact one-sided $p$-value on discordant
pairs (wins vs.\ losses), which directly tests the directional claim
``augmentation helps more often than it hurts.'' This choice is deliberate: our data
contain many ties in regimes where the base oracle is already strong, and tied-heavy
settings can make signed-rank tests less aligned with this directional question.

\paragraph{Shared mechanics.}
All experiments use $C=\mathbb{R}^n$, so no projection is needed. The exploration step
uses $\gamma=1$, $r=1$, accepts a candidate $x+\hat t v$ when
\[
  h(x+\hat t v)+\tfrac{\gamma}{2}\hat t^2 < h(x),
\]
and then greedily keeps the better of the oracle point and accepted exploration point.
The step size is sampled as $\hat t \sim \mathrm{Unif}([0,r])$ in all reported experiments.
The sphere sampler uses $v\sim \mathrm{Unif}(\mathbb{S}^{n-1})$. The gauss-axis sampler
first chooses a coordinate index $i$ uniformly, then samples
$g\sim \mathcal{N}(0, I + (\mu^2-1)e_i e_i^\top)$, and sets $v=g/\|g\|$; at $\mu=1$ this
reduces to sphere sampling, while larger $\mu$ concentrates samples near coordinate axes.
Whenever both samplers are run, we state explicitly which sampler is reported in the
aggregate table.

\subsection{D-stationarity diagnostics}
\label{subsec:exp-dstat-checks}

The counts ``\# non-d-stat'' in \Cref{tab:experiments-story-summary} flag terminal
iterates that fail a numerical stationarity test run once per random instance (on
the final iterate of each run). The three benchmarks use different objectives, so
the stationarity checks differ; each is defined in the following paragraphs.

\paragraph{Trimmed lasso (exact DC inclusion test).}
We use the standard DC decomposition
$h(x)=\tfrac12\|Ax-b\|^2+\lambda\|x\|_1-\lambda\,\mathrm{top}_k(|x|)=f(x)-g(x)$ with
$f,g$ convex. For such a model, $x$ is d-stationary if and only if
$\partial g(x)\subseteq \partial f(x)$ \citep{deOliveira2020}, where $\partial f(x)$ and
$\partial g(x)$ denote the convex subdifferentials of $f$ and $g$ at $x$. This inclusion is
equivalent to a single nonnegative scalar being zero: writing
$\mathrm{dist}_2(u,S)\mathrel{:=}\inf_{v\in S}\|u-v\|_2$ for nonempty closed convex
$S\subseteq\R^d$ and $u\in\R^d$ (in particular $d=1$ for each coordinate projection below),
\begin{equation}\label{eq:tl-dstat-dist-sup}
  \partial g(x)\subseteq \partial f(x)
  \quad\Longleftrightarrow\quad
  \sup_{\xi\in\partial g(x)} \mathrm{dist}_2\!\bigl(\xi,\,\partial f(x)\bigr)=0,
\end{equation}
where the supremum is taken over all subgradients $\xi\in\partial g(x)$ (the set
$\partial g(x)$ is compact for this $g$). The diagnostic below evaluates the right-hand
quantity in~\eqref{eq:tl-dstat-dist-sup} in closed form: every $\xi\in\partial g(x)$ is
determined by a top-$k$ support pattern for the magnitudes $|x_i|$, and the distance to
$\partial f(x)$ decouples across coordinates because $\partial f(x)$ is a Cartesian product
of the one-dimensional sets in~\eqref{eq:tl-partial-f}.

Write the residual gradient and the $k$th largest magnitude threshold as
\begin{equation*}
  r \mathrel{:=} A^\top(Ax-b)\in\R^n,
  \qquad
  \tau \mathrel{:=} |x|_{(k)},
\end{equation*}
where $|x|_{(k)}$ denotes the $k$th order statistic of $(|x_1|,\ldots,|x_n|)$ (so
$\mathrm{top}_k(|x|)=\sum_{j=1}^k |x|_{(j)}$). The convex subdifferential of
$f(x)=\tfrac12\|Ax-b\|^2+\lambda\|x\|_1$ is a Cartesian product of one-dimensional
sets; equivalently, its $i$th coordinate projection is
\begin{equation}\label{eq:tl-partial-f}
  (\partial f(x))_i \;=\;
  \begin{cases}
    \{\, r_i + \lambda\,\mathrm{sgn}(x_i)\,\}, & x_i \neq 0,\\[2pt]
    \bigl[\, r_i-\lambda,\; r_i+\lambda \,\bigr], & x_i = 0.
  \end{cases}
\end{equation}
Fix a tie tolerance $\delta_{\mathrm{tie}}\mathrel{:=}10^{-10}$ and define the
top-$k$ magnitude partition
\begin{equation*}
\begin{aligned}
  S_{\mathrm{high}} &\mathrel{:=} \{\, i : |x_i| > \tau + \delta_{\mathrm{tie}} \,\},\\
  S_{\mathrm{tie}} &\mathrel{:=} \{\, i : \bigl| |x_i|-\tau \bigr| \le \delta_{\mathrm{tie}} \,\},\\
  S_{\mathrm{low}} &\mathrel{:=} \{\, i : |x_i| < \tau - \delta_{\mathrm{tie}} \,\}.
\end{aligned}
\end{equation*}
When $|S_{\mathrm{tie}}|=0$ (equivalently, the top-$k$ magnitude set is uniquely determined
up to the tolerance $\delta_{\mathrm{tie}}$, so there is no nontrivial ambiguity at the
cutoff $\tau$), $\partial g(x)$ is a singleton $\{\xi\}$. The unique $\xi\in\partial g(x)$ has coordinates
$\xi_i=\lambda\,\mathrm{sgn}(x_i)$ if $i$ lies among the $k$ largest magnitudes and
$\xi_i=0$ otherwise. For this $\xi$, the coordinate projections of $\partial f(x)$ in~\eqref{eq:tl-partial-f}
yield $\delta_i \mathrel{:=} \mathrm{dist}_2\!\bigl(\xi_i,(\partial f(x))_i\bigr)$, and
\begin{equation*}
  \sup_{\eta\in\partial g(x)} \mathrm{dist}_2\!\bigl(\eta,\,\partial f(x)\bigr)
  \;=\; \mathrm{dist}_2\!\bigl(\xi,\,\partial f(x)\bigr)
  \;=\; \sqrt{\sum_{i=1}^n \delta_i^2}
  \;=:\; \mathrm{gap}(x),
\end{equation*}
the second equality being the product geometry of~\eqref{eq:tl-partial-f}. We declare
\emph{pass} if $\delta_i\le \varepsilon$ for every $i$, with
\begin{equation*}
  \varepsilon \mathrel{:=} 10^{-6}.
\end{equation*}

When $|S_{\mathrm{tie}}|>0$ and the top-$k$ support is genuinely ambiguous (equivalently,
$|S_{\mathrm{high}}|<k<|S_{\mathrm{high}}|+|S_{\mathrm{tie}}|$), the set $\partial g(x)$ is a
polytope generated by choosing which $n_{\mathrm{pick}}\mathrel{:=}k-|S_{\mathrm{high}}|$
indices inside $S_{\mathrm{tie}}$ join the mandatory indices $S_{\mathrm{high}}$ in the
top-$k$ magnitude set. On $S_{\mathrm{high}}$ and $S_{\mathrm{low}}$ the corresponding
$\xi_i\in\partial g(x)$ are fixed at $\xi_i=\lambda\,\mathrm{sgn}(x_i)$ and $\xi_i=0$,
respectively. In particular, on $S_{\mathrm{low}}$ with $x_i=0$ one has
$\mathrm{dist}_2\bigl(0,(\partial f(x))_i\bigr)=\max\{0,|r_i|-\lambda\}$, whereas on
$S_{\mathrm{low}}$ with $x_i\neq 0$ one has
$\mathrm{dist}_2\bigl(0,(\partial f(x))_i\bigr)=|r_i+\lambda\,\mathrm{sgn}(x_i)|$ because
$(\partial f(x))_i$ is the singleton $\{r_i+\lambda\,\mathrm{sgn}(x_i)\}$ (this is generally
\emph{not} the same as $|r_i|$). On $S_{\mathrm{tie}}$ each admissible vertex uses $p_i\in\{0,1\}$ to indicate
whether coordinate $i$ is included in the top-$k$ support; for each $i\in S_{\mathrm{tie}}$
define the squared coordinate distance if $p_i$ is chosen,
\begin{equation*}
  D_i(p)^2 \mathrel{:=}
  \begin{cases}
    \Bigl(\max\bigl\{\,0,\; |r_i|-\lambda(1-p)\,\bigr\}\Bigr)^2, & x_i = 0,\; p\in\{0,1\},\\[4pt]
    \Bigl(\lambda\,\mathrm{sgn}(x_i)(p-1)-r_i\Bigr)^2, & x_i \neq 0,\; p\in\{0,1\},
  \end{cases}
\end{equation*}
where the lower branch is the squared distance from
$\xi_i=\lambda\,\mathrm{sgn}(x_i)\,p$ to the singleton $(\partial f(x))_i=\{r_i+\lambda\,\mathrm{sgn}(x_i)\}$.
Because $(\partial f(x))_i$ is an interval (or a point) independently across $i$, for any
$\xi\in\partial g(x)$ compatible with tie indicators $(p_i)_{i\in S_{\mathrm{tie}}}$ one has
\begin{equation*}
  \mathrm{dist}_2\!\bigl(\xi,\,\partial f(x)\bigr)^2 \;=\; \sum_{i=1}^n D_i(p_i)^2,
\end{equation*}
with the understanding that $p_i$ is fixed on $S_{\mathrm{high}}\cup S_{\mathrm{low}}$ by the
forced pattern ($p_i=1$ on $S_{\mathrm{high}}$, $p_i=0$ on $S_{\mathrm{low}}$) and only the tied
coordinates vary. Thus maximizing $\mathrm{dist}_2(\xi,\partial f(x))$ over $\xi\in\partial g(x)$
reduces to maximizing this separable squared distance over the tie polytope. The maximizer is
attained at a $0/1$ vertex; it is found by the greedy rule
$M_i\mathrel{:=}D_i(1)^2-D_i(0)^2$, sort $\{M_i\}_{i\in S_{\mathrm{tie}}}$ in decreasing order,
set $p_i=1$ on the first $n_{\mathrm{pick}}$ indices and $p_i=0$ on the remaining tied
indices. With $(p_i)$ fixed, let $\xi\in\partial g(x)$ be the corresponding vertex, set
$\delta_i$ on $S_{\mathrm{tie}}$ to $\sqrt{D_i(p_i)^2}$, combine with the
already-fixed coordinates on $S_{\mathrm{high}}$ and $S_{\mathrm{low}}$, and define
\begin{equation*}
  \mathrm{gap}(x) \mathrel{:=} \sqrt{\sum_{i=1}^n \delta_i^2}
  \;=\; \sup_{\eta\in\partial g(x)} \mathrm{dist}_2\!\bigl(\eta,\,\partial f(x)\bigr).
\end{equation*}
The same pass criterion $\delta_i\le \varepsilon$ coordinatewise with $\varepsilon=10^{-6}$
certifies $\mathrm{gap}(x)=0$ up to floating-point tolerance, hence~\eqref{eq:tl-dstat-dist-sup}.

\paragraph{Least trimmed squares (exact smooth-$f$ DC test).}
Here $h(x)=\tfrac12\|Ax-b\|^2-\tfrac12\,\mathrm{top}_q((Ax-b)^2)=f(x)-g(x)$ with $f$
smooth convex and $g$ convex. In this template, d-stationarity is equivalent to
$g$ being differentiable at $x$ and satisfying $\nabla f(x)=\nabla g(x)$: if $g$ is
not differentiable, $\partial g(x)$ is not a singleton and cannot be contained in
the singleton $\partial f(x)=\{\nabla f(x)\}$. We form squared residuals
$r_i^2$, partition indices into above-, at-, and below-threshold groups relative
to the top-$q$ cutoff, and set $\nabla f(x)=A^\top r$. If there is no residual tie,
the top-$q$ active set is unique, $\nabla g(x)$ is explicit, and the gap is
$\|\nabla f(x)-\nabla g(x)\|_2$. If there are ties, we first test
one admissible top-$q$ completion; if it already violates gradient match beyond
$10^{-6}$, the point is declared non-d-stationary immediately. Otherwise we measure
nondifferentiability by the maximum, over admissible completions $T$ of the tied
block, of $\|\nabla g_T(x)-\nabla g_{T_{\mathrm{ref}}}(x)\|_2$, where
$\nabla g_T(x)=A^\top w$ with $w_i=r_i$ on $T$ and $w_i=0$ off $T$. Exhaustive
enumeration is used when the number of completions is at most $8192$; for larger tie
sets a deterministic pair of extreme completions plus $256$ random completions
(estimating the witness). The acceptance tolerance is $10^{-6}$ on the resulting gap.

\paragraph{ReLU regression (exact DC inclusion check).}
For ReLU regression we now use an explicit DC decomposition and test d-stationarity
through convex-subdifferential inclusion. With
\begin{equation*}
  h(x)=\tfrac12\sum_{i=1}^m\bigl(\max\{0,a_i^\top x\}-b_i\bigr)^2,
\end{equation*}
define $h=f-g$ by
\begin{equation*}
\begin{aligned}
  f(x) &\mathrel{:=} \tfrac12\sum_{i=1}^m\bigl(\max\{0,a_i^\top x\}\bigr)^2
  + \sum_{i=1}^m \alpha_i \max\{0,a_i^\top x\},\\
  g(x) &\mathrel{:=} \sum_{i=1}^m \beta_i \max\{0,a_i^\top x\},
\end{aligned}
\end{equation*}
where $\beta_i \mathrel{:=} b_i$ for $b_i>0$ and $0$ otherwise, and
$\alpha_i \mathrel{:=} -b_i$ for $b_i<0$ and $0$ otherwise. Then
$\alpha_i\beta_i=0$ for all $i$, and d-stationarity is the DC condition
$\partial g(x)\subseteq\partial f(x)$.

Write $z\mathrel{:=}Ax$, and use the numerical kink bandwidth
\begin{equation*}
  \tau_{\mathrm{kink}}(x) \mathrel{:=} 10^{-6}\,\max\bigl\{\,1,\;\|z\|_\infty\,\bigr\}.
\end{equation*}
Set $I_+(x)\mathrel{:=}\{i:z_i>\tau_{\mathrm{kink}}(x)\}$ and
$I_0(x)\mathrel{:=}\{i:|z_i|\le\tau_{\mathrm{kink}}(x)\}$, and split
$I_0(x)=I_{0,\alpha}(x)\sqcup I_{0,\beta}(x)$ with
\begin{equation*}
  I_{0,\alpha}(x)\mathrel{:=}\{i\in I_0(x):\alpha_i>0\},
  \qquad
  I_{0,\beta}(x)\mathrel{:=}\{i\in I_0(x):\beta_i>0\}.
\end{equation*}
Define
\begin{equation*}
\begin{aligned}
  g_0(x) &\mathrel{:=} \sum_{i\in I_+(x)} \beta_i a_i,\\
  f_0(x) &\mathrel{:=} \sum_{i\in I_+(x)} \bigl(z_i+\alpha_i\bigr)a_i,\\
  D(x)   &\mathrel{:=} g_0(x)-f_0(x),
\end{aligned}
\end{equation*}
and generator matrices
\begin{equation*}
  G_\beta(x)\mathrel{:=}\bigl[\beta_i a_i\bigr]_{i\in I_{0,\beta}(x)},
  \qquad
  F_\alpha(x)\mathrel{:=}\bigl[\alpha_i a_i\bigr]_{i\in I_{0,\alpha}(x)}.
\end{equation*}
Then $\partial g(x)\subseteq\partial f(x)$ is equivalent to
\begin{equation*}
  D(x)+G_\beta(x)\lambda \in \bigl\{F_\alpha(x)\mu:\mu\in[0,1]^{|I_{0,\alpha}(x)|}\bigr\}
  \quad \forall \lambda\in[0,1]^{|I_{0,\beta}(x)|}.
\end{equation*}
The implemented diagnostic evaluates the exact vertex form of this condition:
\begin{equation*}
  \Gamma_{\mathrm{ReLU}}(x)\mathrel{:=}
  \max_{\lambda\in\{0,1\}^{|I_{0,\beta}(x)|}}
  \min_{\mu\in[0,1]^{|I_{0,\alpha}(x)|}}
  \bigl\|D(x)+G_\beta(x)\lambda-F_\alpha(x)\mu\bigr\|_2.
\end{equation*}
Each inner problem is a box-constrained least-squares subproblem, and the outer
maximum is over binary vertices. We declare \emph{pass} iff
\begin{equation*}
  \Gamma_{\mathrm{ReLU}}(x)\le 10^{-6}.
\end{equation*}
When $|I_{0,\beta}(x)|$ exceeds the implementation cap (set to $16$), the run is marked
non-d-stationary for reporting rather than using an uncontrolled approximation.

\subsection{Trimmed Lasso: Resolving the Failure Mode}

The trimmed-lasso objective
\[
  h(x)=\tfrac12\|Ax-b\|^2+\lambda\bigl(\|x\|_1-\mathrm{top}_k|x|\bigr)
\]
is the regime where plain DCA is known to converge to critical points that can fail
d-stationarity. We focus on $\lambda=1$ with
$m=50$, $n=100$, $k=5$, noise std.\ $0.1$, and $N=5000$ outer iterations.
Both samplers are run; the aggregate table reports gauss-axis ($\mu=300$).

\begin{center}
% Auto-generated by scripts/make_figures.py. Do not edit by hand.
\begin{tabular}{l c c c c c c}
\toprule
$\lambda$ & $n$ & win/tie/loss & mean gap & median gap (wins) & McNemar $p$ & \# non-d-stat (base/aug) \\
\midrule
1.00 & 100 & 99/1/0 & +2.94e-01 & +2.69e-01 & 1.58e-30 & 99/4 \\
\bottomrule
\end{tabular}

\end{center}

The central empirical message is clear: DCA fails the d-stationarity diagnostic on
$99/100$ instances, while augmentation with gauss-axis reduces this to $4/100$.
At the objective level, augmentation wins on $99$ instances, ties on one, and never
loses. McNemar's exact one-sided test gives $p\approx 1.6\times 10^{-30}$.

Sphere runs were still checked for completeness, but we do not plot
them here because they add little visual information in this focused regime:
the sphere variant accepts no exploration moves and is practically indistinguishable
from DCA. Instead, we show only the gauss-axis plots in
\Cref{fig:v2-tl-scatter,fig:v2-tl-trajectory}.

The reason this happens is structural. In high dimension, a random sphere direction has
typical $\ell_1$ mass of order $\sqrt{n}$ after $\ell_2$ normalization. For trimmed
lasso, the nonconvex sparsity term is driven by coordinate-level behavior, so this
``diffuse'' direction tends to spread motion over many coordinates and rarely aligns with
the sparse descent geometry near stuck points. The acceptance test then almost never
fires in finite budgets. The axis-biased Gaussian mixture is designed to correct exactly
this mismatch: each sampled direction is still absolutely continuous on the sphere and
keeps full support (so the theorem assumptions are unchanged), but it places much larger
probability on directions close to coordinate axes, which are precisely the directions
that can unlock descent for this $\ell_1$-structured objective. Empirically, this
sampler-design change is the difference between inert exploration and systematic
improvement.

\begin{figure}[t]
\centering
\includegraphics[width=0.58\linewidth]{experiments/figures/fig_scatter_gauss_axis.pdf}
\caption{Trimmed lasso scatter at $\lambda=1$. Horizontal axis: final DCA objective.
Vertical axis: median augmented objective over three seeds. Points below the diagonal
favor augmentation. (Sphere sampler omitted from display: in this regime it produces no
accepted exploration moves and overlays DCA.)}
\label{fig:v2-tl-scatter}
\end{figure}

\begin{figure}[t]
\centering
\includegraphics[width=0.62\linewidth]{experiments/figures/fig_trajectory_gauss_axis.pdf}
\caption{Trimmed lasso trajectory on a representative instance. Black: DCA.
Colored band/curve: min/median/max over augmented seeds. The Gaussian-mixture sampler
breaks away from the DCA plateau. (Sphere sampler omitted from display: trajectory
essentially overlays DCA.)}
\label{fig:v2-tl-trajectory}
\end{figure}

\subsection{Least Trimmed Squares: Non-Disruptive Behavior with Occasional Escapes}

For least trimmed squares,
\[
  h(x)=\tfrac12\|Ax-b\|^2-\tfrac12\,\mathrm{top}_q((Ax-b)^2),
\]
we focus on $\sigma_{\mathrm{clean}}=4$ with
$m=100$, $n=50$, $q=10$, outlier std.\ $10$, and $N=5000$.
For the focused LTS report, exploration directions are sampled from the sphere
distribution $v\sim \mathrm{Unif}(\mathbb{S}^{n-1})$.

\begin{center}
% Auto-generated by scripts/make_figures_lts.py.
% Columns: sigma_clean = inlier noise std (not a Tikhonov lambda).
\begin{tabular}{l c c c c c c}
\toprule
$\sigma_{\mathrm{clean}}$ & $n$ & win/tie/loss & mean gap & median gap (wins) & McNemar $p$ & \# non-d-stat (base/aug) \\
\midrule
4.00 & 100 & 9/91/0 & +3.96e+00 & +5.22e+01 & 1.95e-03 & 0/0 \\
\bottomrule
\end{tabular}

\end{center}

This experiment supports a different part of the story. Here the base DCA oracle is
already reliable in first-order terms: both base and augmented runs have zero
d-stationarity failures in the focused setting. Augmentation is therefore expected to be
mostly neutral, and that is exactly what we observe ($91$ ties). Yet on the discordant
subset, all outcomes favor augmentation ($9$ wins, $0$ losses), yielding
$p\approx 2.0\times 10^{-3}$ by McNemar's exact test. In short, augmentation is
non-disruptive where the oracle already behaves well, but can still find a better basin
when one is nearby.

\begin{figure}[t]
\centering
\includegraphics[width=0.45\linewidth]{experiments/figures/fig_lts_scatter.pdf}
\caption{LTS scatter at $\sigma_{\mathrm{clean}}=4$ (sphere reporting).
Most points lie on the diagonal (ties), with a smaller set below (augmentation wins).}
\label{fig:v2-lts-scatter}
\end{figure}

\begin{figure}[t]
\centering
\includegraphics[width=0.6\linewidth]{experiments/figures/fig_lts_trajectory.pdf}
\caption{LTS representative trajectory where augmentation improves over the DCA plateau.}
\label{fig:v2-lts-trajectory}
\end{figure}

\subsection{ReLU Regression: Transfer Beyond DC Oracles}

The ReLU experiment is intentionally outside the DC setup:
\[
  h(x)=\tfrac12\sum_i\bigl(\mathrm{relu}(a_i^\top x)-b_i\bigr)^2,
\]
optimized by a prox-linear oracle ($\rho_{\mathrm{prox}}=0.1$), then augmented by the
same exploration mechanism used elsewhere. We aggregate two focused settings,
$(q,\rho_b)\in\{(0.2,2),(0.4,2)\}$, with $100$ instances per setting
($200$ paired comparisons total), $m=200$, $n=50$, and $N_{\mathrm{outer}}=1000$.
For the focused ReLU report, exploration directions are likewise sampled as
$v\sim \mathrm{Unif}(\mathbb{S}^{n-1})$.

\begin{center}
% Auto-generated by scripts/make_figures_relu_regression.py.
\begin{tabular}{l c c c c c c}
\toprule
setting & $n$ & win/tie/loss & mean gap & median gap (wins) & McNemar $p$ & \# non-d-stat (base/aug) \\
\midrule
focused sweep & 200 & 44/156/0 & +9.58e-01 & +1.26e+00 & 5.68e-14 & 0/0 \\
\bottomrule
\end{tabular}

\end{center}

As in LTS, augmentation is largely neutral but never harmful in the focused report:
$44$ wins, $156$ ties, $0$ losses, with zero d-stationarity failures for both base and
augmented methods. McNemar's exact one-sided test gives
$p\approx 5.7\times 10^{-14}$. Since the inner oracle here is prox-linear (not DCA), this
result is direct empirical evidence for the paper's oracle-agnostic claim.

\begin{figure}[t]
\centering
\includegraphics[width=0.42\linewidth]{experiments/figures/fig_relu_scatter.pdf}
\caption{ReLU scatter (sphere reporting). Most pairs tie, with a smaller strict-win set.}
\label{fig:v2-relu-scatter}
\end{figure}

\begin{figure}[t]
\centering
\includegraphics[width=0.55\linewidth]{experiments/figures/fig_relu_trajectory.pdf}
\caption{ReLU representative trajectory (base prox-linear vs.\ augmented seeds).}
\label{fig:v2-relu-trajectory}
\end{figure}

\subsection{Cross-Experiment Interpretation}

Taken together, the three experiments show a consistent pattern with distinct roles.
Trimmed lasso demonstrates the failure mode and its correction: augmentation converts a
high non-d-stationarity regime into a mostly d-stationary one, with strong objective
improvements. LTS and ReLU demonstrate the complementary regime: when the base oracle is
already first-order sound, augmentation is mostly neutral, never harmful in the focused
runs, and occasionally beneficial.

This is the practical reading of the theoretical separation in the paper.
Oracle choice controls local descent structure and rates; exploration controls asymptotic
d-stationarity behavior in an oracle-agnostic way. The experiments were designed so that
each benchmark contributes one part of that statement, and the combined suite supports it
without requiring any algorithm-specific retuning of the exploration step.

% \input{experiments_appendix}
% \newpage
% \input{neurips_checklist}

\end{document}
